% Format-only conversion of the empirical-extension manuscript at f13d803.
% No LaTeX compilation, PDF generation, rendering, or visual validation performed.
% Figures are byte-identical original SVGs. Future compilation requires the
% standard svg package and Inkscape; authorised shell escape may be required.
% Evidence hyperlinks resolve relative to this directory in a repository checkout.
\documentclass[11pt,a4paper]{article}
\usepackage[T1]{fontenc}
\usepackage[utf8]{inputenc}
\usepackage{lmodern}
\usepackage{textcomp}
\usepackage[margin=25mm]{geometry}
\usepackage{amsmath,amssymb,amsthm}
\usepackage{graphicx,svg}
\usepackage{array,booktabs,longtable,tabularx,xltabular}
\usepackage{caption}
\usepackage{algorithm}
\usepackage{fvextra}
\usepackage{xurl}
\usepackage[hidelinks]{hyperref}
\newtheorem{proposition}{Proposition}
\newcolumntype{Y}{>{\raggedright\arraybackslash}X}
\setcounter{secnumdepth}{2}
\newcommand{\appendixsection}[2]{%
  \refstepcounter{section}%
  \section*{Appendix \thesection. #2}%
  \addcontentsline{toc}{section}{Appendix \thesection. #2}%
  \label{#1}%
  \setcounter{table}{0}%
}
\title{TrafficTwin: Admission and Dispatch Semantics in Vehicular Edge Computing under a Frozen MAPPO Policy}
\author{}
\date{}
\begin{document}
\maketitle
% BEGIN SOURCE BLOCK 002 paragraph
S M Abdulla Al Mamun\\
The University of Manchester\\
Master's dissertation --- empirical-extension revision for author review, 8 September 2026\\
Supervisor: Dr Sandra Sampaio
% END SOURCE BLOCK 002

% BEGIN SOURCE BLOCK 003 paragraph
This revision preserves the manuscript at \texttt{ee6a4b151d92329f9c26e3f344d87b70b84fc31e} and scientific evidence at \texttt{04f3b6a95c7bc06c80ed95b54762f12861bba183}. New type aggregation and scheduler timings are retrospective; the joint-randomness campaign is sealed but unrun. \href{\detokenize{../empirical_extension_2026-09-08/AUTHOR_INPUTS.md}}{Author inputs} identify pending attribution, assessment-specific AI-use and submission decisions. No approval or declaration is implied.
% END SOURCE BLOCK 003

% BEGIN SOURCE BLOCK 004 heading
\section*{Abstract}\phantomsection\label{sec:abstract}
% END SOURCE BLOCK 004

% BEGIN SOURCE BLOCK 005 paragraph
Infrastructure scheduling can reverse a vehicular computing comparison even when the learned vehicle policy is frozen. TrafficTwin studies this through matched strongest-link ingress, common-target least-workload and causal per-task placement. Across four incident fleet draws, common-target falls 2.122 percentage points below gate-enabled ingress, while an adaptive per-task extension exceeds it by 0.527 points. A separate four-draw morning replication excludes its inspected pilot and yields \ensuremath{-}3.232 and +3.899 points, with simultaneous intervals supporting both directions.
% END SOURCE BLOCK 005

% BEGIN SOURCE BLOCK 006 paragraph
Operational analysis distinguishes shared destinations, immediate reservations and admission reconciliation. In the exact studied model, two deadline values guarantee that three simultaneous replacements equal causal fixed-target admission; service and timing constraints explain queue clearing and recurring common destinations. Constructed float32 boundaries qualify executable equivalence without establishing historical prevalence.
% END SOURCE BLOCK 006

% BEGIN SOURCE BLOCK 007 paragraph
Retrospective morning accounting identifies 284,821 gains and 12,842 losses. New task-type analysis locates almost all net gains in the two distinct 100 ms types; one draw also shows increased admitted misses within a benefiting type. These transitions do not identify individual-decision causal effects. A bounded compiled scheduler benchmark finds lower CPU cost for causal per-task scans on the studied fixtures and a further reduction for cyclic targeting. The cyclic comparator is implemented and tested, but its full deadline comparison and fresh joint-randomness confirmation remain unrun.
% END SOURCE BLOCK 007

% BEGIN SOURCE BLOCK 008 paragraph
The contribution is an inspectable implementation study, not a new learned controller. Limited fleet replication, global service-work knowledge and aggregate timing restrict generalisation. The archived capacity contrast remains statistically inconclusive and separately measurement-qualified by an unquantified scoring-mask exception whose original task records are unavailable.
% END SOURCE BLOCK 008

% BEGIN SOURCE BLOCK 009 heading
\section{Introduction}\label{sec:1}
% END SOURCE BLOCK 009

% BEGIN SOURCE BLOCK 010 heading
\subsection{Problem and motivation}\label{sec:1.1}
% END SOURCE BLOCK 010

% BEGIN SOURCE BLOCK 011 paragraph
Vehicular edge computing couples radio and computational resources: a strong radio connection can reach a busy server, and accepting work does not ensure timely completion \cite{ref1}. TrafficTwin asks which implementation decisions cause an observed scheduling advantage.
% END SOURCE BLOCK 011

% BEGIN SOURCE BLOCK 012 paragraph
TrafficTwin combines a trace-driven VEC evaluator with an evidence-oriented software artifact. A supplied MAPPO actor selects Local, V2I or V2V; infrastructure logic selects the execution RSU after radio ingress. This boundary exposes downstream scheduling without retraining the vehicle policy \hyperref[source-s4]{[S4]}.
% END SOURCE BLOCK 012

% BEGIN SOURCE BLOCK 013 paragraph
The initial capacity question distinguished waiting-room space from useful service, motivating E0 accounting and the offered-task denominator. Opposite rankings for common-target and per-task least-workload implementations then made decision and reservation semantics central.
% END SOURCE BLOCK 013

% BEGIN SOURCE BLOCK 014 heading
\subsection{Related work and the position of this study}\label{sec:1.2}
% END SOURCE BLOCK 014

% BEGIN SOURCE BLOCK 015 paragraph
The closest intellectual comparison is task dispatch to parallel servers, particularly the information used to estimate waiting time and the point at which an assignment changes that information. RL supplies the upstream policy, but it does not define the downstream scheduling problem. The following synthesis positions the completed experiments retrospectively: these papers were newly inspected during manuscript revision and are not presented as having guided the historical experiments.
% END SOURCE BLOCK 015

% BEGIN SOURCE BLOCK 016 paragraph
\textbf{Queue count and remaining work.} Shortest-queue dispatch compares numbers of jobs; shortest-workload dispatch compares their outstanding service requirements. With heterogeneous task sizes, these orderings can disagree. Harchol-Balter, Crovella and Murta explicitly study dynamic least-work-remaining assignment when service demand is known, alongside random, round-robin and size-based assignment. Their model assumes immediate per-arrival assignment, identical hosts, non-preemptive first-come-first-served service and Poisson arrivals. Their analysis shows that preferred assignment policies depend on task-size variability \cite{ref2}. Consequently, neither choosing the least workload nor achieving balanced allocation establishes universal optimality. TrafficTwin's inherited \texttt{jsq} name is misleading if read literally: the implemented selector compares service milliseconds, while a separate task count enforces finite admission capacity.
% END SOURCE BLOCK 016

% BEGIN SOURCE BLOCK 017 paragraph
\textbf{Batch decisions and reservations.} Batching does not inherently imply sending a batch to one server. Sparrow pools probes across a parallel job and spreads its tasks over selected workers; late binding queues reservations and assigns tasks when workers become ready. Its discussion identifies both queue-count prediction errors and races between schedulers using apparently idle workers \cite{ref9}. This distinguishes two issues that a generic ``batch versus per-task'' description would conflate: sharing information across candidates can help, whereas committing many assignments against one unchanged destination can concentrate work. TrafficTwin evaluates the latter implementation. Its immediate service-work reservation is also different from Sparrow's worker-triggered task binding: TrafficTwin assumes the necessary information and acknowledgement are already available.
% END SOURCE BLOCK 017

% BEGIN SOURCE BLOCK 018 paragraph
Ying, Srikant and Kang make the reservation distinction especially explicit: batch-filling samples queue counts, assigns tasks one by one and updates the chosen queue after each assignment. Their analysis assumes identical servers, exponential service, Poisson batch arrivals and random ties \cite{ref13}. This is close prior art for sequential within-batch updates. Although the evaluator comments mention their batch-filling work, its common-target reconciliation does not implement that destination rule. TrafficTwin's causal change is therefore an adaptation of established dispatch semantics, evaluated with workload information and deadline admission, rather than a novel batching principle.
% END SOURCE BLOCK 018

% BEGIN SOURCE BLOCK 019 paragraph
\textbf{Local information and common destinations.} Vargaftik, Keslassy and Orda's Local Shortest Queue (LSQ) policies maintain possibly outdated queue-count estimates at multiple dispatchers. Their slotted model can send all jobs arriving at one dispatcher in a slot to one locally shortest queue, with random tie-breaking. Algorithms update local estimates by the jobs sent and refresh selected entries through communication. Their stability result requires stated arrival/service assumptions and bounded expected estimation error \cite{ref10}. This is particularly useful counterevidence to an overbroad novelty claim: common destinations and local updates both predate TrafficTwin. What differs here is service-work information, deterministic index ties, a backlog gate, five ordered task slots and one-second aggregate draining. Neither the LSQ stability theorem nor its distributed communication results transfer directly to these finite-horizon deadline outcomes.
% END SOURCE BLOCK 019

% BEGIN SOURCE BLOCK 020 paragraph
\textbf{Admission and useful completion.} Niño-Mora jointly studies admission and routing to heterogeneous multiserver queues, charging separately for rejection and admitted deadline misses. The model uses Poisson arrivals, exponential service, soft deadlines and state-dependent index policies; admitted late jobs remain until completion \cite{ref11}. This establishes admission as an optimisation decision with consequences beyond the admitted population. TrafficTwin instead preserves a simple inherited backlog filter to make the placement contrast interpretable. The gate does not estimate the complete response time or promise success. Counting successes over all offered tasks makes a reject-all policy score zero and prevents selective admission from concealing failures. That denominator is a deliberate measurement choice, not a newly invented admission algorithm.
% END SOURCE BLOCK 020

% BEGIN SOURCE BLOCK 021 paragraph
Table~\ref{tab:1} locates the comparison along operational dimensions. The papers identify relevant mechanisms and alternatives; they are not external benchmark arms. Their arrival processes, objectives and resource models differ too much for published performance percentages to be compared as controlled results.
% END SOURCE BLOCK 021

% BEGIN SOURCE BLOCK 022 table
\begingroup
\small
\begin{xltabular}{\textwidth}{@{}YYYY@{}}
\caption{Closest-work comparison. ``Work'' means remaining service demand; ``count'' means queued jobs. TrafficTwin rows describe the frozen implementation, not a scheduling-family definition.}\label{tab:1} \\
\toprule
Work / decision layer and arrivals & Information and update semantics & Admission / objective & Relationship to TrafficTwin \\
\midrule\endfirsthead
\toprule
Work / decision layer and arrivals & Information and update semantics & Admission / objective & Relationship to TrafficTwin \\
\midrule\endhead
\bottomrule\endfoot
Harchol-Balter et al. \cite{ref2}: dispatch each arrival & Known service work; central per-arrival choice & Mean waiting time; size-normalised waiting & Closest least-workload comparator; no vehicular gate \\
Sparrow \cite{ref9}: parallel-job task placement & Sampled workers; queued reservations; worker replies bind tasks & Job response time; placement constraints & Batch spreading differs from one common destination; communication is explicit \\
Ying et al. \cite{ref13}: tasks within an arriving batch & Sampled counts; update after each assignment; random ties & Delay and sampling cost & Closest within-batch update rule; no deadline gate \\
LSQ \cite{ref10}: dispatcher batch per time slot & Local queue counts; own-assignment increments and sampled/pushed updates & Stability under specified assumptions & Common routing is established; random ties and distributed estimates differ \\
Niño-Mora \cite{ref11}: joint admission/routing per arrival & Queue counts and model-based indices & Rejection and deadline-miss costs & Supports joint evaluation, not the inherited backlog predicate \\
TrafficTwin common-target: one choice per substep & Global service work; vectorised candidate offsets and three reconciliations & Backlog gate plus finite task limit & Frozen vehicle policy; instantaneous logical forwarding \\
TrafficTwin per-task: ordered candidate scan & Global service work; actual admitted work reserved immediately & Same predicate and limit, causal admission & Tested implementation change; global information is assumed \\
\end{xltabular}
\endgroup
% END SOURCE BLOCK 022

% BEGIN SOURCE BLOCK 023 paragraph
\textbf{Attribution and simulation validity.} Sargent distinguishes verification of a computer model from operational validation for its intended application \cite{ref12}. TrafficTwin's conservation and replay checks primarily establish the former; observed physical RSU service and network measurements would be needed for stronger deployment claims. SUMO supplies a microscopic simulation framework \cite{ref8}, but importing its trajectories does not independently validate a coupled radio/computing model. This distinction explains why a source-correct result can be scientifically useful while remaining conditional on simulation semantics.
% END SOURCE BLOCK 023

% BEGIN SOURCE BLOCK 024 paragraph
PPO and MAPPO identify the reused actor's provenance \cite{ref3}, \cite{ref4}, without establishing checkpoint optimality or identical actions under changed observations. Henderson et al., Agarwal et al. and Gorsane et al. motivate exposing implementation choices and finite-run uncertainty \cite{ref5}, \cite{ref6}, \cite{ref7}. TrafficTwin therefore checks matched actions and uses fleet draws as replications; those references do not authorise its particular small-sample distributional assumption.
% END SOURCE BLOCK 024

% BEGIN SOURCE BLOCK 025 paragraph
The contribution concerns a defined control boundary: established workload and reservation ideas yield opposite rankings against ingress. Decision semantics and offered-task accounting make that result interpretable.
% END SOURCE BLOCK 025

% BEGIN SOURCE BLOCK 026 heading
\subsection{Aim, research questions and contribution}\label{sec:1.3}
% END SOURCE BLOCK 026

% BEGIN SOURCE BLOCK 027 paragraph
The aim is to determine whether deterministic infrastructure-side RSU load management improves offered-task deadline attainment under a frozen vehicle policy, and to establish which claims the available evidence can support. The questions developed with the programme; the morning replication questions and analysis were declared before its new outcomes, while the mechanism audit was explicitly post-hoc.
% END SOURCE BLOCK 027

% BEGIN SOURCE BLOCK 028 paragraph
\textbf{RQ1 --- Measurement and capacity.} What do the recorded accounting checks establish, and does increasing the RSU waiting-room limit at fixed compute service improve the archived score? Interpretation as admission-consistent deadline attainment additionally depends on the unresolved legacy scoring-mask exception (\S\,\ref{sec:3.1}). Sampling and measurement uncertainty are separate questions.
% END SOURCE BLOCK 028

% BEGIN SOURCE BLOCK 029 paragraph
\textbf{RQ2 --- Admission and dispatch semantics.} Under matched inputs and the inherited workload-based admission rule, how do ingress execution, common-target least-busy placement and sequential per-task least-busy placement compare in the incident scenario? The distinction concerns the implemented decision process, including its reconciliation semantics.
% END SOURCE BLOCK 029

% BEGIN SOURCE BLOCK 030 paragraph
\textbf{RQ3 --- Replication and mechanism.} Does the incident ordering recur across newly declared fleet draws in the morning scenario, and what recorded or reconstructable behaviour explains the common-target scheduler's repeatedly selected RSU identities? The replication and the subsequent explanatory audit have different evidential roles.
% END SOURCE BLOCK 030

% BEGIN SOURCE BLOCK 031 paragraph
\textbf{RQ4 --- Bounded sensitivities.} Within one incident fleet draw, how do aged placement-workload reports and a fixed forwarding overhead affect the comparison? These studies test specific model assumptions descriptively, without establishing population-level robustness or a physical backhaul design.
% END SOURCE BLOCK 031

% BEGIN SOURCE BLOCK 032 paragraph
The contributions are source-audited accounting; adaptive incident comparisons and separate morning replication; and an exact-model mechanism explanation with numerical limits. Current retrospective extensions add task-type interpretation, measured scheduler costs and a tested cyclic comparator with an unrun joint-randomness protocol. They ask whether workload awareness adds value beyond spreading, and whether reversal persists under fresh joint randomness; neither question is answered by existing full-study results. Personal attribution is separately assessed \hyperref[source-s11]{[S11]}, \hyperref[source-s13]{[S13]}--\hyperref[source-s15]{[S15]}.
% END SOURCE BLOCK 032

% BEGIN SOURCE BLOCK 033 heading
\subsection{Scope and report structure}\label{sec:1.4}
% END SOURCE BLOCK 033

% BEGIN SOURCE BLOCK 034 paragraph
One actor, a provisional fleet model and two Manchester traces bound the study. Section~\ref{sec:2} specifies the model and design; Section~\ref{sec:3} evaluates findings and achievement; Section~\ref{sec:4} answers the questions. Appendices support verification.
% END SOURCE BLOCK 034

% BEGIN SOURCE BLOCK 035 heading
\section{Methodology}\label{sec:2}
% END SOURCE BLOCK 035

% BEGIN SOURCE BLOCK 036 heading
\subsection{Research design and evidence hierarchy}\label{sec:2.1}
% END SOURCE BLOCK 036

% BEGIN SOURCE BLOCK 037 paragraph
Matched traffic, fleet and task inputs isolate infrastructure changes under a frozen policy; recorded actions are checked rather than assumed identical. Conclusions concern this evaluator, not retrained or physically validated systems.
% END SOURCE BLOCK 037

% BEGIN SOURCE BLOCK 038 paragraph
The sequence is adaptive: earlier observations informed later questions. E2c excludes its discovery seed, while the morning replication excludes its inspected pilot. Frozen source and accepted records determine experimental meaning; correspondence clarifies intention. The completed raw audit and subsequent kernel analyses are retrospective investigations, not part of the original protocols \hyperref[source-s0]{[S0]}--\hyperref[source-s8]{[S8]}, \hyperref[source-s13]{[S13]}.
% END SOURCE BLOCK 038

% BEGIN SOURCE BLOCK 039 heading
\subsection{Architecture and information boundaries}\label{sec:2.2}
% END SOURCE BLOCK 039

% BEGIN SOURCE BLOCK 040 paragraph
Figure~\ref{fig:1} separates actor mode choice, environmental radio targeting and infrastructure execution/admission. V2I ingress is the strongest simulated radio link. Ingress execution retains that RSU; load-balancing arms may assign elsewhere. A target remains a proposal until admission \hyperref[source-s4]{[S4]}.
% END SOURCE BLOCK 040

% BEGIN SOURCE BLOCK 041 figure
\begin{figure}[htbp]
\centering
\includesvg[width=\linewidth]{assets/architecture}
\caption{Experimental responsibility boundaries. Only admitted V2I work enters the selected execution queue; logical forwarding applies when ingress and execution differ. Resource scaling is outside the completed comparison.}\label{fig:1}
\end{figure}
% END SOURCE BLOCK 041

% BEGIN SOURCE BLOCK 042 paragraph
The 17-dimensional observation includes task descriptors, vehicle queues, state of charge, radio summaries, compute capability, EV status and observed V2V target tier, but excludes RSU workload. One mode per active vehicle-second serves all independently sampled operational tasks in its five slots. Admission can nevertheless affect transmit energy, then SoC and later observations/actions. Frozen weights therefore do not guarantee fixed actions; the old matched-action finding is empirical, and the new protocol allows this feedback.
% END SOURCE BLOCK 042

% BEGIN SOURCE BLOCK 043 paragraph
V2V target selection is environmental: exclude the source and peers whose queues are full, then select the remaining strongest instantaneous simulated link. Distance contributes to link quality, but fading also contributes; this is not simply nearest-neighbour selection. Peer workload and capability affect latency after selection, while capability is not itself the target-ranking criterion. Observation and operational link calculations use separate fading samples. The correspondence confirms this and the per-vehicle-second decision convention as intended inherited behaviour, not an accidental change introduced by the placement study \hyperref[source-s4]{[S4]}.
% END SOURCE BLOCK 043

% BEGIN SOURCE BLOCK 044 heading
\subsection{Scenarios, actor and controlled inputs}\label{sec:2.3}
% END SOURCE BLOCK 044

% BEGIN SOURCE BLOCK 045 paragraph
Table~\ref{tab:2} identifies the two scenarios. Traffic positions are imported from saved traces; the evaluator does not rerun SUMO during each scheduling comparison. The morning input audit reconstructs the canonical trace arrays from their source records and checks entry markers and slot assignment. This matters because padded array slots are not necessarily permanent vehicle identities. On morning vehicle entry, the recorded convention resets vehicle queues for the new visit; state of charge is not reset by that queue operation. The incident trace retains its earlier convention without the same entry channel \hyperref[source-s7]{[S7]}.
% END SOURCE BLOCK 045

% BEGIN SOURCE BLOCK 046 table
\begingroup
\small
\begin{xltabular}{\textwidth}{@{}YYY@{}}
\caption{Scenario identity and fixed placement controls. Differences between scenarios are retained and documented, not treated as a single-factor intervention.}\label{tab:2} \\
\toprule
Setting & Incident scenario & Morning scenario \\
\midrule\endfirsthead
\toprule
Setting & Incident scenario & Morning scenario \\
\midrule\endhead
\bottomrule\endfoot
Date and local window & 15 March 2024, 20:00--21:00 & 15 October 2024, 08:00--11:00 \\
One-second steps & 3,600 & 10,800 \\
Padded vehicle slots & 2,488 & 215 \\
RSUs & 10 & 9 \\
SUMO seed & 43 & 42 \\
Vehicle queue convention & Conserved, legacy trace convention & Conserved, canonical per-visit reset \\
Placement-study RSU limit & 6,220 tasks per RSU & 6,220 tasks per RSU \\
Service / forwarding / scaling & 1\ensuremath{\times} / 0 ms / off & 1\ensuremath{\times} / 0 ms / off \\
\end{xltabular}
\endgroup
% END SOURCE BLOCK 046

% BEGIN SOURCE BLOCK 047 paragraph
The same one-hot-17 MAPPO checkpoint, identified by its archived hash, is retained throughout. Its inherited training provenance is Model C with synthetic mobility, 128 environments, learning rate 0.003 and training seed 100. Those facts identify the input artifact; they are not a training reproduction performed here. The provisional UK2030 fleet preset samples RPi-, Jetson- and GPU-vehicle tiers with probabilities (0.40, 0.35, 0.25), and electric-vehicle status with probability 0.22. These are modelling assumptions, not measured Manchester fleet shares. Fleet seeds vary between replicates, while evaluator seed 0 remains fixed. Thus the estimand concerns fleet variability conditional on the actor and scenario, not uncertainty over training or traffic generation.
% END SOURCE BLOCK 047

% BEGIN SOURCE BLOCK 048 paragraph
Each active vehicle offers $\min(\operatorname{Poisson}(1.5),5)$ tasks per second; inactive slots offer none. Operational task types are independently sampled with probabilities (0.20, 0.30, 0.50). Types 1--3 have deadlines (100, 500, 100) ms, mean input sizes (1, 0.0012, 0.001) MB and workloads (1,254, 2,100, 15) million cycles. Input sizes multiply their type mean by an independent uniform 0.8--1.2 factor. The five-slot cap changes offered work as well as scheduling opportunities \hyperref[source-s4]{[S4]}.
% END SOURCE BLOCK 048

% BEGIN SOURCE BLOCK 049 paragraph
For homogeneous RSUs, service milliseconds are $s = \frac{\xi \times C_t}{2.2 \times 1.5 \times \min(12,p_t) \times 0.9 \times \min(5,g_t)}$, where $C_t$ is million cycles, $p_t=(4,4,2)$, $g_t=(40,5,1)$, and \ensuremath{\xi} is uniform on 0.9--1.1. The denominator uses 2.2 GHz, 1.5 instructions per cycle, effective cores, utilisation and capped GPU acceleration. Million-cycle/GHz conversion yields milliseconds. Before noise, services are approximately (21.111, 35.354, 2.525) ms. All placement arms use service multiplier 1.
% END SOURCE BLOCK 049

% BEGIN SOURCE BLOCK 050 paragraph
Absolute capacity avoids a confound: the incident ratio 2.5 \ensuremath{\times} 2,488 gives 6,220 tasks per RSU; applying 2.5 to the morning width would give only 538. The morning protocol fixes 6,220 explicitly. Density, duration, layout and vehicle-entry conventions still differ.
% END SOURCE BLOCK 050

% BEGIN SOURCE BLOCK 051 heading
\subsection{Time, queues, admission and outcomes}\label{sec:2.4}
% END SOURCE BLOCK 051

% BEGIN SOURCE BLOCK 052 paragraph
A batch advances simulated time by one second. Five internal task substeps order admissions; they are not five 200 ms clock advances. Work accumulates across them before one service drain. Let $W[r]$ denote remaining service milliseconds at RSU $r$, $Q[r]$ its aggregate task count and C the task limit. These variables serve different purposes: identical counts can represent different work, and increasing C does not accelerate service \hyperref[source-s4]{[S4]}.
% END SOURCE BLOCK 052

% BEGIN SOURCE BLOCK 053 paragraph
For each offered task $i$, let $a[i]$ indicate admission, $y[i]$ indicate the recorded deadline success, and $z[i,c]$ indicate terminal failure category c. The intended accounting contract requires one admission or failure and no success for rejected work. The per-run primary outcome and validation requirements are:
% END SOURCE BLOCK 053

% BEGIN SOURCE BLOCK 054 equation
\begin{equation}\label{eq:1}
 A=\frac{N_{\mathrm{met}}}{N_{\mathrm{offered}}}
  =\frac{\sum_i y[i]}{N_{\mathrm{offered}}};
 \qquad \text{reported percentage}=100A.
\end{equation}
% END SOURCE BLOCK 054

% BEGIN SOURCE BLOCK 055 equation
\begin{equation}\label{eq:2}
 a[i]+\sum_c z[i,c]=1;\qquad y[i]\leq a[i];\qquad
 N_{\mathrm{offered}}=N_{\mathrm{admitted}}+N_{\mathrm{failed}}.
\end{equation}
% END SOURCE BLOCK 055

% BEGIN SOURCE BLOCK 056 paragraph
Here $N_{\mathrm{failed}}$ includes both rejection and unavailability. Categories distinguish local queue rejection, V2V queue rejection or unavailability, and V2I unavailability, backlog-gate rejection or capacity rejection. Gate failure takes precedence over capacity failure when both apply. Selected targets are proposals; rejected tasks neither execute nor reserve work (Figure~\ref{fig:2}). Admitted attainment uses $N_{\mathrm{admitted}}$ instead, so must not replace Equation~\eqref{eq:1}. Offered, admitted and successful-task latency means similarly have different populations.
% END SOURCE BLOCK 056

% BEGIN SOURCE BLOCK 057 figure
\begin{figure}[htbp]
\centering
\includesvg[width=\linewidth]{assets/task_accounting}
\caption{Offered tasks remain in the primary denominator. Admission, execution assignment and deadline outcome are separate fields; selecting an RSU does not establish that work executed there.}\label{fig:2}
\end{figure}
% END SOURCE BLOCK 057

% BEGIN SOURCE BLOCK 058 paragraph
The inherited gate tests existing effective workload, including the candidate offset $O[i]$, against deadline $D[i]$:
% END SOURCE BLOCK 058

% BEGIN SOURCE BLOCK 059 equation
\begin{equation}\label{eq:3}
 W[r]+O[i]<D[i],\quad\text{together with an available task-count place.}
\end{equation}
% END SOURCE BLOCK 059

% BEGIN SOURCE BLOCK 060 paragraph
The inequality is strict. It excludes the candidate's service, radio transfer and forwarding cost, and therefore is a backlog filter rather than an end-to-end feasibility guarantee. Candidate ordering and offsets differ between algorithms (\S\,\ref{sec:2.5}). Admission uses positive ingress link quality as its coarse radio condition; the latency model additionally requires quality at least 0.15 for a usable transmission. Consequently, admission can still lead to a transmission-related deadline miss.
% END SOURCE BLOCK 060

% BEGIN SOURCE BLOCK 061 paragraph
For viable V2I execution, the modelled response latency is:
% END SOURCE BLOCK 061

% BEGIN SOURCE BLOCK 062 equation
\begin{equation}\label{eq:4}
 \begin{aligned}
 L_{\mathrm{raw}}[i]={}&T(\mathrm{input}[i],q,c)+W[r]+O[i]+s[i,r]\\
 &+T(0.001,q,c)+f\cdot\mathbf{1}(r\neq\mathrm{ingress}[i]).
 \end{aligned}
\end{equation}
% END SOURCE BLOCK 062

% BEGIN SOURCE BLOCK 063 paragraph
$T(b,q,c)$ is $0.1 + 1{,}000 \times 8b/c$ milliseconds for $q\geq 0.15$ and capacity $c>0.000000001$ Mbps, with b in MB; otherwise it uses a $10^{9}$ ms failure value (up to the 0.1 ms propagation addition). The 0.001 MB return term reuses the ingress quality and capacity. It is not a separately simulated downlink or inter-RSU route. The fixed forwarding charge f is added once. Local latency is vehicle backlog plus local service; V2V adds input/return transfers on the selected peer link, peer backlog, preceding local/V2V reservations and peer service. Service $s$ derives from task cycles, device frequency, parallelism, utilisation and acceleration, with a sampled 0.9--1.1 multiplier. The placement code reserves the same sampled work used for latency: exact service knowledge is a declared information assumption.
% END SOURCE BLOCK 063

% BEGIN SOURCE BLOCK 064 paragraph
The gate-enabled paths pass final admission into latency scoring. Historical \texttt{off} instead passes earlier coarse eligibility while enqueueing uses refined admission: an in-substep capacity rejection need not receive penalty latency. Its historical numerical impact remains unquantified (\S\,\ref{sec:3.1}). The final mapping is $L[i]=10D[i]$ when $L_{\mathrm{raw}}[i]\geq 10^{8}$ ms, and $L[i]=L_{\mathrm{raw}}[i]$ otherwise. Finite late latencies are not generally clipped; success uses $L[i]\leq D[i]$, inclusively. Finite rejection latency alone is not a false deadline success.
% END SOURCE BLOCK 064

% BEGIN SOURCE BLOCK 065 paragraph
For each queue over an interval, work accounting requires:
% END SOURCE BLOCK 065

% BEGIN SOURCE BLOCK 066 equation
\begin{equation}\label{eq:5}
 W_{\mathrm{initial}}+\sum\text{admitted service}
 =\sum\text{served service}+W_{\mathrm{final}}.
\end{equation}
% END SOURCE BLOCK 066

% BEGIN SOURCE BLOCK 067 paragraph
After the fifth substep, $S[r]=\min(W_{\mathrm{post}}[r],1{,}000)$ is served and $W_{\mathrm{next}}[r]=W_{\mathrm{post}}[r]-S[r]$. If this empties the queue, its task count becomes zero. Otherwise the evaluator subtracts $\left\lfloor Q_{\mathrm{post}}[r]\times S[r]/W_{\mathrm{post}}[r]\right\rfloor$, retaining at least one task while positive work remains. This aggregate carry approximation has no individual departure-event model. Equation~\eqref{eq:5} checks the declared dynamics, not their physical calibration.
% END SOURCE BLOCK 067

% BEGIN SOURCE BLOCK 068 heading
\subsection{Placement algorithms and their implementation}\label{sec:2.5}
% END SOURCE BLOCK 068

% BEGIN SOURCE BLOCK 069 paragraph
Table~\ref{tab:3} separates execution destination from the gate switch. The identifier \texttt{jsq} actually selects least service workload, not shortest task count. \texttt{dla} combines this common-target selector with Equation~\eqref{eq:3}; these inherited names do not define canonical scheduling algorithms \hyperref[source-s2]{[S2]}, \hyperref[source-s4]{[S4]}.
% END SOURCE BLOCK 069

% BEGIN SOURCE BLOCK 070 table
\begingroup
\small
\begin{xltabular}{\textwidth}{@{}YYYY@{}}
\caption{Infrastructure conditions. All retain strongest-link radio ingress. ``Live'' eligibility is refreshed at each substep; the historical off mode retains step-entry coarse eligibility.}\label{tab:3} \\
\toprule
Mode & Execution destination & Backlog gate & Coarse eligibility \\
\midrule\endfirsthead
\toprule
Mode & Execution destination & Backlog gate & Coarse eligibility \\
\midrule\endhead
\bottomrule\endfoot
\texttt{off} & Ingress & Off & Step entry \\
\texttt{jsq} & Common least-workload target & Off & Live \\
\texttt{ingress\_dla} & Ingress & On & Live \\
\texttt{dla} & Common least-workload target & On & Live \\
\texttt{per\_task\_dla} & Causal least-workload target & On & Live \\
\end{xltabular}
\endgroup
% END SOURCE BLOCK 070

% BEGIN SOURCE BLOCK 071 paragraph
Both placement algorithms below operate once per task substep, using its entry state $W$, Q. Index $i$ follows ascending padded-vehicle order, not arrival-time or deadline priority. Each candidate has actual service $s[i,r]$, deadline $D[i]$ and a fixed actor action and radio ingress. Every argmin considers all RSUs and breaks ties by lowest index; a full selected RSU is rejected rather than excluded and retried elsewhere.
% END SOURCE BLOCK 071

% BEGIN SOURCE BLOCK 072 paragraph
For common-target reconciliation, let $E[i]$ mean active V2I attempt, positive ingress quality and $Q[r[i]]<C$ at substep entry. For an admission mask $M$, define $O[i,M]$ as the sum of $s[j,r[i]]$ over earlier $j<i$ with the same target and $M[j]$ true; $H[i,M]$ is their count. Offsets are exclusive: a candidate's own service is omitted. Algorithm~\ref{alg:1} follows the frozen implementation's three simultaneous replacements, without claiming convergence.
% END SOURCE BLOCK 072

% BEGIN SOURCE BLOCK 073 algorithm
\begin{algorithm}[htbp]
\caption{common-target placement with gate (one substep)}\label{alg:1}
\begin{Verbatim}[fontsize=\small,breaklines=true]
r[i] = lowest-index argmin(W), shared by all candidates
M = E
repeat exactly three times:
    compute O[i,M] and H[i,M] from the previous whole mask
    M_new[i] = E[i] and (W[r[i]] + O[i,M] < D[i])
                       and (H[i,M] < C - Q[r[i]])
    M = M_new simultaneously
recompute final offsets O[i,M] for latency
label E-candidates failing the final workload test as gate failures
label remaining E-candidates excluded by M as capacity failures
commit each M-admitted task and its full service once
unavailable or rejected tasks reserve nothing
\end{Verbatim}
\end{algorithm}
% END SOURCE BLOCK 073

% BEGIN SOURCE BLOCK 074 paragraph
Target selection never changes during reconciliation. The final admission mask and the offsets recomputed from it determine the recorded processing; three passes are an implementation constant. Section~\ref{sec:3.4} establishes the exact two-deadline guarantee and its float32 boundary. Ingress with the gate uses the same offset/reconciliation machinery with individual ingress targets instead of the common argmin. Gate-off modes omit the workload test.
% END SOURCE BLOCK 074

% BEGIN SOURCE BLOCK 075 algorithm
\begin{algorithm}[htbp]
\caption{causal per-task placement with gate (one substep)}\label{alg:2}
\begin{Verbatim}[fontsize=\small,breaklines=true]
repeat exactly three complete passes:
    working_W = W; working_Q = Q       # restart, not cumulative
    for i in ascending padded-vehicle order:
        r[i] = lowest-index argmin(working_W)
        O[i] = working_W[r[i]] - W[r[i]]
        require active V2I, positive quality and Q[r[i]] < C
        test working_W[r[i]] < D[i], then working_Q[r[i]] < C
        if admitted:
            working_W[r[i]] += s[i,r[i]]
            working_Q[r[i]] += 1
        otherwise label failure; reserve nothing
retain the last pass; commit its admitted work once
\end{Verbatim}
\end{algorithm}
% END SOURCE BLOCK 075

% BEGIN SOURCE BLOCK 076 paragraph
The causal gate and capacity see earlier admissions immediately. Repeated passes begin from identical initial state and reproduce the same decisions; they do not triple reservations. Physical queue arrays receive the final result once, before the next substep. Both algorithms drain service only after all five substeps. Thus the comparison changes selection frequency, reservation visibility and reconciliation semantics; it is not an isolated change to argmin frequency.
% END SOURCE BLOCK 076

% BEGIN SOURCE BLOCK 077 paragraph
For an illustrative scheduling calculation, assume three empty identical RSUs, four viable 40 ms tasks with 100 ms deadlines, nonbinding capacity and negligible communication. Common-target chooses RSU 0; its masks settle at three admissions and one gate rejection. The admitted offsets are 0, 40 and 80 ms, leaving 120, 0, 0 ms: the third admitted task misses its deadline despite passing the backlog gate. Per-task selects 0, 1, 2, 0, leaves 80, 40, 40 ms and all four meet deadlines under these assumptions. Figure~\ref{fig:3} depicts the proposals. This constructed example explains mechanics; it is neither a recorded task group nor an estimate of recoverable historical failures.
% END SOURCE BLOCK 077

% BEGIN SOURCE BLOCK 078 figure
\begin{figure}[htbp]
\centering
\includesvg[width=\linewidth]{assets/dispatch_example}
\caption{Four illustrative proposals. The common destination remains fixed through reconciliation; causal placement updates after admissions. The accompanying worked example states admissions and deadline outcomes under deliberately simplified assumptions.}\label{fig:3}
\end{figure}
% END SOURCE BLOCK 078

% BEGIN SOURCE BLOCK 079 heading
\subsection{Experimental progression and inference}\label{sec:2.6}
% END SOURCE BLOCK 079

% BEGIN SOURCE BLOCK 080 paragraph
E0 validates recorded accounting before performance analysis. E1 compares three queue limits at fixed service across five matched fleet draws. E2 and E2b are exploratory single-draw studies separating placement and admission. E2c evaluates common-target against ingress on four new incident draws, seeds 1--4. E2d is an adaptive extension after inspecting E2c: it adds per-task placement on those previously examined draws and reuses the exact controls. They are not new independent observations \hyperref[source-s0]{[S0]}--\hyperref[source-s3]{[S3]}.
% END SOURCE BLOCK 080

% BEGIN SOURCE BLOCK 081 paragraph
The morning pilot uses seed 1 for ingress and per-task placement. After inspecting it, the replication declares seeds 0, 2, 3 and 4 and all three arms, yielding twelve new full runs. The primary contrasts are per-task minus ingress, common-target minus ingress, and per-task minus common-target. The pilot is excluded; a supplementary two-arm five-draw summary is descriptive. The local protocol and manifest preceded the new outcomes, but there was no external preregistration \hyperref[source-s7]{[S7]}, \hyperref[source-s8]{[S8]}.
% END SOURCE BLOCK 081

% BEGIN SOURCE BLOCK 082 paragraph
For each contrast, a fleet draw supplies one paired difference in percentage points. Draws receive equal weight. The mean difference has a Student-t interval, $\text{mean}\pm t\times\frac{\text{sample standard deviation}}{\sqrt{n}}$. Incident E2c/E2d retain their reported individual 95\% intervals. The morning protocol additionally uses Bonferroni simultaneous intervals for three contrasts, with critical value t at probability $1-0.05/(2\times 3)$, three degrees of freedom. Its reversal criterion requires the simultaneous common-target-minus-ingress interval below zero and the per-task-minus-ingress interval above zero. Four draws cannot strongly diagnose distributional assumptions; these small-sample intervals remain conditional parametric summaries, not task-level significance tests.
% END SOURCE BLOCK 082

% BEGIN SOURCE BLOCK 083 heading
\subsection{Sensitivities, mechanism audit and validation}\label{sec:2.7}
% END SOURCE BLOCK 083

% BEGIN SOURCE BLOCK 084 paragraph
The state-information pilot uses incident seed 1 and report ages 0, 100, 500 and 1,000 ms, alongside ingress. Admission remains live. For previous post-admission workload B and positive age $d$, the report is $\max(B-(1{,}000-d),0)$, with current-batch reservations added immediately. This assumes continuous service between batches without redistributing arrivals. Startup uses live state and recorded age zero \hyperref[source-s5]{[S5]}.
% END SOURCE BLOCK 084

% BEGIN SOURCE BLOCK 085 paragraph
Four ten-step forwarding probes qualify fixed-latency transformations of saved full records: charges leave admission, queues and actions unchanged. No full positive-cost simulation or network model is implied; penalty-equal ambiguity remains unresolved \hyperref[source-s6]{[S6]}.
% END SOURCE BLOCK 085

% BEGIN SOURCE BLOCK 086 paragraph
The reused gap-closure audit authenticated 87 September arrays and checked 19 full runs: twelve morning primary, two excluded pilot and five incident sensitivity cells. Active masks define offers; enqueue fields define V2I admission independently of success. Non-V2I admission excludes rejection predicates. Flags, categories, latency/deadlines, assignments and summaries are cross-checked; execution fields record admission, not departures \hyperref[source-s13]{[S13]}.
% END SOURCE BLOCK 086

% BEGIN SOURCE BLOCK 087 paragraph
The prior joins require identical manifest/fleet and trace-row/substep/slot coordinates, with matching task, arrival, action and ingress fields. Unsaved sizes correspond through the arm-independent random-key schedule. Gains minus losses must equal the archived success difference over a common offered denominator; this is descriptive accounting.
% END SOURCE BLOCK 087

% BEGIN SOURCE BLOCK 088 paragraph
The existing reference distinguishes simultaneous common-target admission A, causal admission replaying A's targets B, and causal reselection C. The completed exact proof and 719-input numerical investigation are reused; Appendix~\ref{app:C} records their scope. Masks require exact agreement, independently of offset tolerance. No prior audit, task join or numerical suite is repeated \hyperref[source-s14]{[S14]}.
% END SOURCE BLOCK 088

% BEGIN SOURCE BLOCK 089 heading
\subsection{Retrospective extensions and measurement design}\label{sec:2.8}
% END SOURCE BLOCK 089

% BEGIN SOURCE BLOCK 090 paragraph
The new causal round-robin comparator isolates workload-aware selection from simple spreading. A pointer starts at RSU 0 and persists across candidates, substeps and batches. Each active V2I attempt with positive ingress quality proposes the pointer's RSU and advances it modulo $R$, even if capacity or the strict gate subsequently rejects it. Unavailable radio and padding do not advance it; there is no retry or full-node skipping. Candidate order, sampled work, gate/capacity rules, precision, accounting, forwarding and drain match causal per-task placement. Repeated passes restart from the same pointer and queues; the final pointer commits once. The new versioned copy preserves frozen sources. Kernel/route checks cover rejection, radio, capacity and carry; full evaluator qualification remains unrun \hyperref[source-s15]{[S15]}.
% END SOURCE BLOCK 090

% BEGIN SOURCE BLOCK 091 paragraph
The scheduler benchmark uses actual JAX helpers, including the unchanged three-replacement prefix operations. Precomputed candidate work replaces only their service provider. Sixteen configurations combine four arms, $N/R=215/9$ or 2,488/10, $K=5$, and sparse/empty or busy/nonempty fixtures. Types retain coupled deadlines and service ranges; all arms receive identical arrays. These are constructed inputs, not observed trajectories. Busy batch-entry queues are imposed stress states, excluded by exact clearing from empty initialisation (\S\,\ref{sec:3.4}). Following official JAX guidance, compilation is separate, inputs are device-resident, five warm-ups precede 100 timed calls, and every output is synchronised. Actor, radio, service generation, transfers, scoring and I/O are excluded. Timing repetitions measure host variability, not fleet uncertainty \cite{ref14}.
% END SOURCE BLOCK 091

% BEGIN SOURCE BLOCK 092 paragraph
New type aggregation authenticates eight September primary-morning task files, one at a time. Outcome categories supply admission, gate rejection and admitted misses; their independent admission validation comes from the prior audit. Type totals must reproduce all-task summaries and paired net changes. All three types and four draws are retained, without new joins or subgroup significance tests.
% END SOURCE BLOCK 092

% BEGIN SOURCE BLOCK 093 heading
\section{Evaluation and Reflection}\label{sec:3}
% END SOURCE BLOCK 093

% BEGIN SOURCE BLOCK 094 heading
\subsection{Accounting validity and the capacity question}\label{sec:3.1}
% END SOURCE BLOCK 094

% BEGIN SOURCE BLOCK 095 paragraph
E0 addressed a concrete mismatch between task scoring and queue admission. In the legacy clamp path, the evaluator scored eligible V2I tasks and multiplied their combined service work by the fraction fitting the task-count ceiling. Surplus tasks were not individually identified in earlier scoring. The 5 August repair introduced per-task admission/rejection and full enqueueing of admitted work. Source history credits Randy's implementation and Abdulla's motivating accounting analysis; E0 qualified recorded accounting in that configuration \hyperref[source-s0]{[S0]}, \hyperref[source-s12]{[S12]}.
% END SOURCE BLOCK 095

% BEGIN SOURCE BLOCK 096 paragraph
An illustrative reconstruction shows why the distinction matters. Suppose one queue place remains and two otherwise viable candidates each require 40 ms. Legacy fractional enqueue adds (1/2) \ensuremath{\times} 80 = 40 ms, yet both candidates can remain in the previously scored population. The queue represents only half their combined demand. Corrected reject-mode admission accepts one candidate in the declared order, rejects the other, enqueues the accepted 40 ms in full and gives the rejected task a failure outcome. These are explanatory numbers, not a historical measured task pair. With unequal service times, fractional enqueue can additionally differ from the work of the particular task actually admitted.
% END SOURCE BLOCK 096

% BEGIN SOURCE BLOCK 097 paragraph
The relevant invariant is that offered tasks partition into admitted and declared failed attempts, and rejected tasks do not execute. E0's validator cross-checks active records, terminal categories and summary counts. Its archived regression fixture changes offered count from two to three without adding a record and requires failure; another changes offered vehicle work from 10 to 11 ms against an 8 + 2 ms partition. However, summary work partition alone is weak evidence of dynamics because rejected work is calculated as offered minus admitted. Later independent reconstruction checks enqueue, drain and endpoints directly.
% END SOURCE BLOCK 097

% BEGIN SOURCE BLOCK 098 paragraph
E0's archived receipts report passing checks, including deadline-flag/outcome agreement. Source versions identify the remaining exception: E0/E1 use \texttt{0f01f4d}, E2 uses \texttt{e11f444}, and E2b reuses E2's off cell. E0 and E2 off each record 834,120 capacity rejections, establishing exposure rather than false successes. Author-reported deletion of original E0/E1/E2 outputs, with no known backup, makes fresh task checks and rescoring not assessable. Surviving receipts cannot quantify non-penalty rejection latency \hyperref[source-s13]{[S13]}.
% END SOURCE BLOCK 098

% BEGIN SOURCE BLOCK 099 paragraph
E1 varied waiting-room capacity across five matched draws. The archived score difference for 99,520 versus 1,866 tasks per RSU was \ensuremath{-}0.01094 percentage points, with a 95\% interval [\ensuremath{-}0.02465, +0.00276]. Recalculation from surviving run summaries reproduces this interval; it is not fresh task-level validation. The archived E1 score contrast is statistically inconclusive under the historical implementation. Its interpretation as admission-consistent deadline attainment remains qualified by the unquantified scoring-mask exception. The interval describes sampling uncertainty conditional on that implementation, not the additional measurement uncertainty. Neither equivalence nor statistically supported degradation follows \hyperref[source-s1]{[S1]}.
% END SOURCE BLOCK 099

% BEGIN SOURCE BLOCK 100 paragraph
The saved secondary summaries report more admissions and lower rejection at larger limits, while mean admitted latency rises from approximately 3.77 to 12.19 and 132.49 seconds. These changing admitted populations cannot substitute for the offered-task contrast; their latency interpretation also retains the mask qualification. Larger waiting rooms can retain work without increasing useful service.
% END SOURCE BLOCK 100

% BEGIN SOURCE BLOCK 101 paragraph
Changing the off scoring mask could also change V2I transmission energy, state of charge and later actor observations. Even an admission-consistent adjustment of saved flags would describe fixed history, not a corrected rerun. Separately, the completed September gate-enabled audit found zero non-admitted successes, zero non-penalty rejection latencies and complete admission/category/summary agreement across its 19 authenticated runs. Their admission-consistent scores equal the archived scores. This finite result supports those later outcomes, without resolving the missing historical records or proving universal correctness.
% END SOURCE BLOCK 101

% BEGIN SOURCE BLOCK 102 heading
\subsection{From admission decomposition to the incident reversal}\label{sec:3.2}
% END SOURCE BLOCK 102

% BEGIN SOURCE BLOCK 103 paragraph
E2's exploratory seed-0 comparison showed offered attainment of approximately 68.362\% for strongest-link execution without the deadline gate, 67.568\% for inherited least-busy placement without the gate, and 69.494\% for inherited least-busy placement with the gate. The last comparison changed both placement and admission relative to the strongest-link reference. Its improvement could not establish that load balancing itself helped \hyperref[source-s2]{[S2]}.
% END SOURCE BLOCK 103

% BEGIN SOURCE BLOCK 104 paragraph
E2b added ingress execution with the same deadline gate, reaching approximately 71.577\%. The strongest-link contrast was +3.215 points and the common-target contrast +1.926 points, giving an interaction of \ensuremath{-}1.290 points. However, \texttt{off} uses step-entry coarse saturation eligibility while \texttt{ingress\_dla} refreshes eligibility per substep. Their contrast and interaction bundle the gate with this timing difference and retain the legacy scoring-mask exception (\S\,\ref{sec:2.4}). The manifest records 138 unavailable V2I attempts among 13,076,234 offered tasks in off; that count does not bound downstream effects. E2b identifies a stronger comparator without assigning its entire gain solely to the gate. Subsequent gate-held placement comparisons use live eligibility and refined scoring masks in all arms. In the exploratory draw, common-target remained 2.083 points below ingress with the gate.
% END SOURCE BLOCK 104

% BEGIN SOURCE BLOCK 105 paragraph
E2c tested the common-target comparison on four new incident fleet draws, excluding discovery seed 0. Its mean difference from ingress was \ensuremath{-}2.122 points, with an individual 95\% interval of [\ensuremath{-}2.234, \ensuremath{-}2.011]. All four paired differences were negative. Inspection then established that the inherited scheduler chose one common target per task substep. E2d introduced the causal per-task implementation under the same actor, trace, gate predicate, queue ceiling, service and forwarding controls. Table~\ref{tab:4} retains all four draw-level results \hyperref[source-s3]{[S3]}.
% END SOURCE BLOCK 105

% BEGIN SOURCE BLOCK 106 table
\begingroup
\small
\begin{xltabular}{\textwidth}{@{}YYYYY@{}}
\caption{Incident offered-task deadline attainment, expressed as percentages. E2d reused the exact E2c ingress and common-target controls; these are four paired draws, not separate sets of independent controls.}\label{tab:4} \\
\toprule
Fleet seed & Ingress & Common-target & Per-task & Per-task \ensuremath{-} ingress, pp \\
\midrule\endfirsthead
\toprule
Fleet seed & Ingress & Common-target & Per-task & Per-task \ensuremath{-} ingress, pp \\
\midrule\endhead
\bottomrule\endfoot
1 & 72.003 & 69.794 & 72.467 & +0.464 \\
2 & 70.369 & 68.317 & 70.956 & +0.587 \\
3 & 70.298 & 68.153 & 70.805 & +0.507 \\
4 & 70.887 & 68.805 & 71.438 & +0.551 \\
\end{xltabular}
\endgroup
% END SOURCE BLOCK 106

% BEGIN SOURCE BLOCK 107 paragraph
Per-task exceeds ingress by 0.527 points, individual 95\% interval [+0.442, +0.612], and common-target by 2.649 points, interval [+2.621, +2.678]. RQ2 therefore yields opposite rankings for the two implementations. These intervals summarise an adaptive extension on examined draws; the morning protocol supplies prospective replication.
% END SOURCE BLOCK 107

% BEGIN SOURCE BLOCK 108 paragraph
Ingress determines the practical scale: against 70.889\% mean attainment, +0.527 points means about 53 extra successes per 10,000 offered, or a 0.744\% relative increase. The larger +2.649-point common-target contrast diagnoses that implementation; it is not the gain over the stronger reference. These are descriptive conversions of the same mean, and neither arm establishes optimal dispatch.
% END SOURCE BLOCK 108

% BEGIN SOURCE BLOCK 109 heading
\subsection{Morning replication with the inspected pilot excluded}\label{sec:3.3}
% END SOURCE BLOCK 109

% BEGIN SOURCE BLOCK 110 paragraph
The inspected morning pilot offered 1,744,116 tasks and produced 91.075\% ingress versus 94.028\% per-task attainment: +2.953 points or 51,503 successes. It motivated replication but remains exploratory \hyperref[source-s7]{[S7]}.
% END SOURCE BLOCK 110

% BEGIN SOURCE BLOCK 111 paragraph
The morning protocol declared four new seeds and all three arms before outcomes. Each draw reproduces common-target < ingress < per-task. Mean attainment is 85.655\%, 88.887\% and 92.785\%, respectively; its lower ingress/per-task means than the pilot underline the need for replication \hyperref[source-s8]{[S8]}.
% END SOURCE BLOCK 111

% BEGIN SOURCE BLOCK 112 table
\begingroup
\small
\begin{xltabular}{\textwidth}{@{}YYYYY@{}}
\caption{Morning primary replication. Seed 1 is absent because it was the inspected pilot; each row is a matched fleet draw on the same morning trace.}\label{tab:5} \\
\toprule
Fleet seed & Ingress, \% & Common-target, \% & Per-task, \% & Per-task \ensuremath{-} ingress, pp \\
\midrule\endfirsthead
\toprule
Fleet seed & Ingress, \% & Common-target, \% & Per-task, \% & Per-task \ensuremath{-} ingress, pp \\
\midrule\endhead
\bottomrule\endfoot
0 & 88.817 & 85.660 & 92.705 & +3.888 \\
2 & 87.536 & 83.537 & 92.060 & +4.524 \\
3 & 89.082 & 85.862 & 92.978 & +3.896 \\
4 & 90.113 & 87.559 & 93.399 & +3.286 \\
\end{xltabular}
\endgroup
% END SOURCE BLOCK 112

% BEGIN SOURCE BLOCK 113 paragraph
Table~\ref{tab:6}'s simultaneous per-task-minus-ingress interval is positive and common-target-minus-ingress interval negative, satisfying the prespecified reversal criterion. Figure~\ref{fig:4} displays these primary results.
% END SOURCE BLOCK 113

% BEGIN SOURCE BLOCK 114 table
\begingroup
\small
\begin{xltabular}{\textwidth}{@{}YYYY@{}}
\caption{Morning paired contrasts in percentage points. Individual intervals are two-sided 95\% Student-t intervals; simultaneous intervals use the prespecified three-contrast Bonferroni adjustment. Replication uses four fleet draws.}\label{tab:6} \\
\toprule
Contrast & Mean & Individual 95\% interval & Simultaneous 95\% interval \\
\midrule\endfirsthead
\toprule
Contrast & Mean & Individual 95\% interval & Simultaneous 95\% interval \\
\midrule\endhead
\bottomrule\endfoot
Per-task \ensuremath{-} ingress & +3.899 & [+3.094, +4.703] & [+2.671, +5.126] \\
Common-target \ensuremath{-} ingress & \ensuremath{-}3.232 & [\ensuremath{-}4.176, \ensuremath{-}2.289] & [\ensuremath{-}4.672, \ensuremath{-}1.793] \\
Per-task \ensuremath{-} common-target & +7.131 & [+5.385, +8.877] & [+4.466, +9.795] \\
\end{xltabular}
\endgroup
% END SOURCE BLOCK 114

% BEGIN SOURCE BLOCK 115 figure
\begin{figure}[htbp]
\centering
\includesvg[width=\linewidth]{assets/primary_replication}
\caption{Redrawn from archived morning primary data. Left: within-draw attainment. Right: draw differences, mean diamonds and three-contrast Bonferroni intervals. The inspected pilot is excluded; uncertainty concerns four newly declared fleet draws, conditional on the frozen actor and morning scenario \hyperref[source-s8]{[S8]}.}\label{fig:4}
\end{figure}
% END SOURCE BLOCK 115

% BEGIN SOURCE BLOCK 116 paragraph
The supplementary five-draw ingress/per-task mean, including the pilot, is +3.709 points descriptively; it does not enter Table~\ref{tab:6}. No common-target pilot result is available.
% END SOURCE BLOCK 116

% BEGIN SOURCE BLOCK 117 paragraph
The morning gain is about 390 extra successes per 10,000 offered, a 4.386\% relative increase over ingress. Its larger magnitude cannot be attributed solely to density: RSU layout/count, date, duration, padded fleet width and entry conventions also change. Separate analyses support recurrence in two Manchester scenarios, without estimating a pooled geographical effect.
% END SOURCE BLOCK 117

% BEGIN SOURCE BLOCK 118 paragraph
Table~\ref{tab:7} makes both sides of the improvement visible over the same 1,744,116 offered tasks per draw. Gains total 284,821 and losses 12,842, giving 271,979 net additional successes. Previously rejected work contributes most gains, but admitted misses also matter; per-task placement introduces smaller losses of both kinds. These are descriptive transitions under different later queues, not individual-decision causal effects. The four fleet draws remain the replications \hyperref[source-s13]{[S13]}, \hyperref[source-s14]{[S14]}.
% END SOURCE BLOCK 118

% BEGIN SOURCE BLOCK 119 table
\begingroup
\small
\begin{xltabular}{\textwidth}{@{}YYYYYY@{}}
\caption{Morning ingress\ensuremath{\to}per-task outcome accounting, re-extracted from the authenticated paired results. Gains become successes; losses cease to be successes. All four primary draws are retained; the pilot is excluded.}\label{tab:7} \\
\toprule
Seed & Gain: gate rejection & Gain: admitted miss & Loss: non-admitted & Loss: admitted miss & Net successes \\
\midrule\endfirsthead
\toprule
Seed & Gain: gate rejection & Gain: admitted miss & Loss: non-admitted & Loss: admitted miss & Net successes \\
\midrule\endhead
\bottomrule\endfoot
0 & 53,061 & 18,091 & 395 & 2,945 & 67,812 \\
2 & 64,418 & 19,125 & 496 & 4,144 & 78,903 \\
3 & 52,432 & 18,430 & 105 & 2,808 & 67,949 \\
4 & 41,823 & 17,441 & 48 & 1,901 & 57,315 \\
Total & 211,734 & 73,087 & 1,044 & 11,798 & 271,979 \\
\end{xltabular}
\endgroup
% END SOURCE BLOCK 119

% BEGIN SOURCE BLOCK 120 table
\begingroup
\small
\begin{xltabular}{\textwidth}{@{}YYYYYY@{}}
\caption{Retrospective type-level net additional successes, per-task minus ingress. Each type has the stated offered population in every primary draw; full successes, admissions and failure counts are supplied in Appendix~\ref{app:D}.}\label{tab:8} \\
\toprule
Type / deadline & Offered per draw & Seed 0 & Seed 2 & Seed 3 & Seed 4 \\
\midrule\endfirsthead
\toprule
Type / deadline & Offered per draw & Seed 0 & Seed 2 & Seed 3 & Seed 4 \\
\midrule\endhead
\bottomrule\endfoot
1 / 100 ms & 349,045 & +19,823 & +19,165 & +20,339 & +20,254 \\
2 / 500 ms & 523,060 & +72 & +62 & +6 & +5 \\
3 / 100 ms & 872,011 & +47,917 & +59,676 & +47,604 & +37,056 \\
\end{xltabular}
\endgroup
% END SOURCE BLOCK 120

% BEGIN SOURCE BLOCK 121 paragraph
The gain is concentrated in the two distinct 100 ms types (Table~\ref{tab:8}). Type 3 supplies 192,253 net successes (70.7\% of the total), Type 1 79,581 (29.3\%), and Type 2 only 145. Within-type gains span 4.249--6.843, 5.491--5.827 and 0.001--0.014 percentage points, respectively. Type 1's much larger payload and service demand distinguish it from Type 3 (\S\,\ref{sec:2.3}). No type has a negative net change in any draw, but this does not mean no tasks are harmed: Table~\ref{tab:7} retains losses. In draw 2, Type 1 gate rejections fall by 23,152 while admitted misses rise by 3,987, yielding 19,165 net successes. Reduced rejection is consistent with access to spare queues; these records cannot isolate deadline length from type-specific communication/service demand or identify each scheduling decision's causal effect \hyperref[source-s15]{[S15]}.
% END SOURCE BLOCK 121

% BEGIN SOURCE BLOCK 122 heading
\subsection{Destination concentration and production-domain mechanisms}\label{sec:3.4}
% END SOURCE BLOCK 122

% BEGIN SOURCE BLOCK 123 paragraph
Workload diagnostics revealed a striking morning contrast. Common-target admitted work used only RSUs 0--4, while per-task placement used all nine. Mean common-target workload shares were approximately 48.34\%, 30.33\%, 14.63\%, 5.21\% and 1.48\% across those five RSUs, with zero work at RSUs 5--8. Per-task shares were close to one ninth each. Figure~\ref{fig:5} is descriptive evidence of allocation, not proof that workload equality alone caused the attainment difference \hyperref[source-s8]{[S8]}.
% END SOURCE BLOCK 123

% BEGIN SOURCE BLOCK 124 figure
\begin{figure}[htbp]
\centering
\includesvg[width=\linewidth]{assets/rsu_workload_distribution}
\caption{Admitted service-work distribution, redrawn from archived CSVs. Bars are means; whiskers are observed minima/maxima over four fleet draws, not confidence intervals. Aggregate shares motivated a separate post-hoc audit; they do not by themselves establish intermediate scheduling decisions or counterfactual task outcomes.}\label{fig:5}
\end{figure}
% END SOURCE BLOCK 124

% BEGIN SOURCE BLOCK 125 paragraph
The completed audit covers four common-target runs: 43,200 seconds and 216,000 substeps. All saved end-of-second workloads and task counts were zero, covering 388,800 entries per field. Every observable target selected RSU k at substep k in 179,427 substeps. The remaining 36,573 had no recorded V2I attempt and do not expose the selector's internal output. Five RSUs formed the full-run union; all five were used in only 15,803 seconds \hyperref[source-s9]{[S9]}.
% END SOURCE BLOCK 125

% BEGIN SOURCE BLOCK 126 paragraph
Intermediate substep workloads were not directly saved for common-target. They were reconstructed from recorded admissions and execution destinations using the previously qualified service-work reference and frozen replay. Every saved endpoint matched exactly, but all-zero endpoints have limited power to distinguish intermediate states. The stronger combined support is the qualified service reference, recorded admissions, frozen enqueue semantics and agreement with every observable target's lowest-index workload minimum. Maximum reconstructed pre-drain workloads ranged from approximately 534.281 to 538.730 ms, below the 1,000 ms service interval; maximum task counts were 22--25, far below 6,220.
% END SOURCE BLOCK 126

% BEGIN SOURCE BLOCK 127 paragraph
A conditional deduction from Algorithm~\ref{alg:1} explains the concentration. With $K$ substeps and $R$ RSUs, at most $\min(K,R)$ distinct destinations receive new assignments in one batch: each substep contributes at most one target, so their union has at most $K$ members. This bound does not require empty queues or a particular tie rule, and concerns new assignments rather than all work already present.
% END SOURCE BLOCK 127

% BEGIN SOURCE BLOCK 128 paragraph
The repeated low-index prefix requires additional assumptions. At an empty batch start, lowest-index ties select RSU 0. Positive admitted work makes that RSU busier than untouched empty nodes; without intervening draining, the next productive substep selects the next empty index. Induction gives a prefix when $R\geq K$; a substep adding no work does not advance it. Complete end-of-batch clearing restores the initial tie. The recorded endpoints and reconstructed maxima satisfy those conditions here. Without clearing or deterministic ties, the per-batch bound could still hold while the full-run union covered all nine RSUs. This is analysis of the stated model, not a universal theorem against batching.
% END SOURCE BLOCK 128

% BEGIN SOURCE BLOCK 129 paragraph
The arrays recorded 511,684 deadline-gate rejections and zero RSU-capacity rejections. In reconstructed states, every gate rejection coexisted with at least five idle RSUs possessing spare task capacity. Larger waiting rooms would not address that gate bottleneck. Spare capacity does not establish that each task would succeed elsewhere: redistribution changes later queues, and end-to-end latency also matters.
% END SOURCE BLOCK 129

% BEGIN SOURCE BLOCK 130 proposition
\begin{proposition}[fixed-target admission in the studied exact model]\label{prop:1}
Fix candidate order, eligibility, destinations, initial nonnegative workload/count and nonnegative service demands, with strict backlog and exclusive count-capacity tests and no intervening drain. At each nonempty destination, let $n$ be the eligible candidate count and $d$ the number of distinct deadlines. Starting from the eligibility mask, simultaneous replacement reaches the unique causal fixed point within $\min(n,2d-1)$ replacements; destinations decouple.\end{proposition}
% END SOURCE BLOCK 130

% BEGIN SOURCE BLOCK 131 proof
\begin{proof}[Proof sketch]
Each candidate depends only on earlier candidates, so a causal scan constructs the unique fixed point and one additional prefix position settles per replacement. The map reverses inclusion. Starting from an upper mask containing the solution, its first extra admission is removed next. Capacity exhaustion ends later admissions; otherwise gate rejection permanently excludes later deadlines no larger than its threshold. After two replacements the unresolved upper-bounded suffix has at most $d-1$ classes. Induction, with one class requiring one replacement, gives the sharper bound. Appendix~\ref{app:C} supplies the complete assumptions/proof \hyperref[source-s14]{[S14]}.\end{proof}
% END SOURCE BLOCK 131

% BEGIN SOURCE BLOCK 132 paragraph
Here ``production'' means the studied evaluator configuration, not deployed roadside hardware. Its two deadline values imply three replacements suffice in exact arithmetic, including nonempty queues and capacity limits. The earlier three-threshold instability is outside this domain. This rules out unrestricted reconciliation failure as an exact-model explanation, while leaving target reselection and executable rounding distinct.
% END SOURCE BLOCK 132

% BEGIN SOURCE BLOCK 133 paragraph
The same model makes clearing structural from empty initialisation. The last admitted task sees less than 500 ms of backlog and contributes at most approximately 38.889 ms of service. Each queue therefore remains below 538.889 ms before a 1,000 ms drain. The observed zero endpoints and reconstructed maxima near 538 ms agree with this consequence. Queue clearing and the resulting destination prefix thus follow from the stated gate/service/timing assumptions, not traffic density alone. This deduction is conditional on exact admission arithmetic and service multiplier 1.
% END SOURCE BLOCK 133

% BEGIN SOURCE BLOCK 134 paragraph
Float32 qualifies this result: inclusive cumulative sums minus current work can round differently from causal exclusive accumulation. The completed investigation found stable third masks and causal-float32 agreement in 324 coupled fixtures; two still differed from exact arithmetic. Four of 210 deliberate threshold fixtures differed from causal float32, with one alternating mask. Imposed-target padding preserved that alternation at studied widths; a common-target-consistent empty-entry fixture remained stable. The alternating task was capacity-labelled despite spare capacity but missed its deadline either way. Sub-0.001 ms errors can change strict gates. These constructed boundaries establish neither full-trajectory reachability nor historical prevalence (Appendix~\ref{app:C}).
% END SOURCE BLOCK 134

% BEGIN SOURCE BLOCK 135 paragraph
Target reselection can separately lose useful completions. Table~\ref{tab:9} uses two empty RSUs, two substeps, capacity 6,220, zero transfer latency and repeated type2/type2/type1/type2 tasks with nominal production service demands. Both arms admit all eight tasks. B replays A's targets 0 then 1, whereas C chooses the current least workload after every reservation.
% END SOURCE BLOCK 135

% BEGIN SOURCE BLOCK 136 table
\begingroup
\small
\begin{xltabular}{\textwidth}{@{}YYYYYYY@{}}
\caption{Reduced-infrastructure illustration with production task/deadline/capacity values. Every listed task is admitted in both arms; completion is diagnostic service latency in milliseconds. ``Slot.task'' preserves candidate order. Values shown to three decimals; tests retain full precision.}\label{tab:9} \\
\toprule
Slot.task & Deadline & Service & B: RSU / work before & B: completion & C: RSU / work before & C: completion \\
\midrule\endfirsthead
\toprule
Slot.task & Deadline & Service & B: RSU / work before & B: completion & C: RSU / work before & C: completion \\
\midrule\endhead
\bottomrule\endfoot
1.1 & 500 & 35.354 & 0 / 0.000 & 35.354 met & 0 / 0.000 & 35.354 met \\
1.2 & 500 & 35.354 & 0 / 35.354 & 70.707 met & 1 / 0.000 & 35.354 met \\
1.3 & 100 & 21.111 & 0 / 70.707 & 91.818 met & 0 / 35.354 & 56.465 met \\
1.4 & 500 & 35.354 & 0 / 91.818 & 127.172 met & 1 / 35.354 & 70.707 met \\
2.1 & 500 & 35.354 & 1 / 0.000 & 35.354 met & 0 / 56.465 & 91.818 met \\
2.2 & 500 & 35.354 & 1 / 35.354 & 70.707 met & 1 / 70.707 & 106.061 met \\
2.3 & 100 & 21.111 & 1 / 70.707 & 91.818 met & 0 / 91.818 & 112.929 \textbf{miss} \\
2.4 & 500 & 35.354 & 1 / 91.818 & 127.172 met & 1 / 106.061 & 141.414 met \\
\end{xltabular}
\endgroup
% END SOURCE BLOCK 136

% BEGIN SOURCE BLOCK 137 paragraph
B preserves an empty RSU for substep 2. C spreads earlier work across both, making the later 100 ms task finish at 112.929 ms: eight admissions produce seven successes. Current workload omits future deadlines, and the gate omits own service; neither directly maximises deadline-success count. Tested deletions remove this loss, without establishing global minimality.
% END SOURCE BLOCK 137

% BEGIN SOURCE BLOCK 138 paragraph
The loss did not persist in twelve tested transfers to nine/ten RSUs and five substeps, with four/six candidates. The example identifies an order/service/deadline interaction, not the cause of recorded losses at full dimensions. B replays A targets and is not an autonomous deployable comparator. Original full-evaluator prefixes remain unrun \hyperref[source-s10]{[S10]}, \hyperref[source-s14]{[S14]}.
% END SOURCE BLOCK 138

% BEGIN SOURCE BLOCK 139 heading
\subsection{Bounded state-information and forwarding sensitivities}\label{sec:3.5}
% END SOURCE BLOCK 139

% BEGIN SOURCE BLOCK 140 paragraph
In incident seed 1, fresh per-task attainment was 72.46695\%, versus ingress 72.00328\%. Report ages 100, 500 and 1,000 ms yielded 72.46695\%, 72.46856\% and 72.45510\%; single-draw changes do not establish equivalence or population robustness \hyperref[source-s5]{[S5]}.
% END SOURCE BLOCK 140

% BEGIN SOURCE BLOCK 141 paragraph
Across 35,990 non-startup RSU/batch observations, live entry workload was zero and every 100 ms report equalled fresh: queues had already emptied. The 100 ms information treatment was inactive. Reports differed at 500 ms in 79.58\% of observations and at 1,000 ms in all. Live admission and immediate acknowledged reservations protect this older-report model; stale capacity, delayed acknowledgements or dispersed arrivals would change it. Equal outcomes therefore cannot establish harmless communication delay.
% END SOURCE BLOCK 141

% BEGIN SOURCE BLOCK 142 paragraph
Fixed forwarding overheads 0, 1, 2.5, 5 and 10 ms yielded 72.46695\%, 72.44783\%, 72.41882\%, 72.37240\% and 72.28579\% attainment. At 10 ms the advantage remained 0.28251 points, or 36,942 successes over ingress, despite 23,689 extra misses relative to zero overhead \hyperref[source-s6]{[S6]}.
% END SOURCE BLOCK 142

% BEGIN SOURCE BLOCK 143 paragraph
These are saved-record transformations qualified by four short direct probes. Forwarding charges do not feed back into admission, queues or actions in that path; the calculation models neither congestion nor postponed arrivals. A further 104 forwarded tasks had penalty-equal latency ambiguity but were already misses and remained misses under every nonnegative charge. Their latency ambiguity remains unresolved. RQ4 supports the tested information and fixed-cost models, not an extrapolated break-even threshold or network design.
% END SOURCE BLOCK 143

% BEGIN SOURCE BLOCK 144 heading
\subsection{Engineering cost, achievement and attribution}\label{sec:3.6}
% END SOURCE BLOCK 144

% BEGIN SOURCE BLOCK 145 table
\begingroup
\small
\begin{xltabular}{\textwidth}{@{}YYYYY@{}}
\caption{Scheduler-only engineering trade-off on Apple M5 CPU, JAX/JAXlib 0.4.30, float32, default runtime threads. Times are milliseconds per five-substep batch; $S$/B denote sparse/busy fixtures. Temporary memory is XLA buffer analysis, not peak RSS.}\label{tab:10} \\
\toprule
Policy & $N=215$ median $S$/B & $N=2{,}488$ median $S$/B & Busy p95, 215/2,488 & Temporary KiB, 215/2,488 \\
\midrule\endfirsthead
\toprule
Policy & $N=215$ median $S$/B & $N=2{,}488$ median $S$/B & Busy p95, 215/2,488 & Temporary KiB, 215/2,488 \\
\midrule\endhead
\bottomrule\endfoot
Ingress & 1.031/1.019 & 6.096/6.113 & 1.103/6.230 & 24.6/294.9 \\
Common-target & 0.952/0.962 & 5.797/5.815 & 1.014/5.935 & 23.8/285.2 \\
Per-task & 0.052/0.055 & 0.501/0.616 & 0.062/0.641 & 11.7/139.1 \\
Round-robin & 0.025/0.021 & 0.223/0.186 & 0.023/0.201 & 11.7/139.1 \\
\end{xltabular}
\endgroup
% END SOURCE BLOCK 145

% BEGIN SOURCE BLOCK 146 paragraph
Table~\ref{tab:10} measures a computational trade-off previously left qualitative. Per-task was faster than the vectorised three-pass helper on these CPU fixtures, while cyclic selection reduced cost further. This does not establish a round-robin deadline advantage: no full comparison has run. Sparse padding still traverses all $N$ positions. Fixed configuration order and uncontrolled host frequency/thermal state limit timing comparisons; raw durations and IQRs are retained.
% END SOURCE BLOCK 146

% BEGIN SOURCE BLOCK 147 paragraph
The vectorised helper materialises $N\times R$ candidate/prefix arrays and performs three replacements: $O(KNR)$ work and $O(NR)$ intermediate space for fixed iteration count, with parallel prefix dependence. Per-task scans have $O(KNR)$ work from $R$-way minima and a $K\times N$ sequential dependency chain; round-robin removes those minima, giving $O(KN)$ target-selection and predicate work in an indexed-state model. Actual queue-update cost depends on compiler buffer reuse. Both keep $R$ queue entries and emit $O(KN)$ diagnostic fields. These are source-level costs; compilers can eliminate unused service broadcasts and repeated identical causal passes. Measured execution therefore need not follow a simple ``vectorised is faster'' ordering.
% END SOURCE BLOCK 147

% BEGIN SOURCE BLOCK 148 paragraph
Compilation took 0.058--0.267 seconds per configuration, separately from lowering and execution. Transfer and process peak memory were not measured. The benchmark excludes full-evaluator costs, and host time is never added to historical simulated task latency. No roadside response-time guarantee follows \hyperref[source-s15]{[S15]}.
% END SOURCE BLOCK 148

% BEGIN SOURCE BLOCK 149 paragraph
Engineering also preserves measurement meaning. Separate proposal, admission and execution fields prevent rejected targets being counted as execution, while reuse of the service subkey avoids resampling a different workload \hyperref[source-s4]{[S4]}.
% END SOURCE BLOCK 149

% BEGIN SOURCE BLOCK 150 paragraph
Table~\ref{tab:11} distinguishes documented intellectual work, supplied code and project implementation. Randy's source explicitly credits Abdulla's accounting questions and clamp reconstruction. The August correspondence records Abdulla's analysis of control boundaries, interpretation of the inconclusive capacity result and explanation of the placement reversal. Those are substantive activities beyond coding; the record does not establish that they were unaided or independently checked. Personal/AI allocation remains for confirmation \hyperref[source-s12]{[S12]}.
% END SOURCE BLOCK 150

% BEGIN SOURCE BLOCK 151 table
\begingroup
\small
\begin{xltabular}{\textwidth}{@{}YYYYY@{}}
\caption{Component ownership and evidence. Attribution certainty concerns what the records establish, not a declaration of independent authorship.}\label{tab:11} \\
\toprule
Component & Inherited starting point & Candidate contribution (project record) & Validation evidence & Attribution certainty \\
\midrule\endfirsthead
\toprule
Component & Inherited starting point & Candidate contribution (project record) & Validation evidence & Attribution certainty \\
\midrule\endhead
\bottomrule\endfoot
Actor, environment and radio targets & Randy's trained actor and VEC model; PPO/MAPPO/JAX reused & Analyse control boundary; hold actor and inputs fixed & Randy's confirmations; actor/input/action checks & Supplied origin clear; no training contribution claimed \\
Accounting corrections & Randy's reject-mode, co-batch, vehicle-queue and lifecycle repairs & Accounting questions and clamp reconstruction credited to Abdulla; project E0 qualification & Commits 4cb7c06/0f01f4d; E0 records and negative fixtures & Requirements credit explicit; repair code supplied, unaided analysis unverified \\
Experimental design and analysis & Existing evaluator and traces & Capacity, admission/placement decomposition, paired replication and bounded follow-ups & Frozen manifests, exclusions, tables and correspondence & Project-level work clear; personal/AI allocation requires confirmation \\
Per-task selector & Common-target selector and three-pass evaluator & Causal selector, matching service reservations, compatibility validation & vec\_env 2f63706; selector tests and production probes & Project change verified; Git author does not prove unaided coding \\
TrafficTwin evidence software & Python, typed models, Streamlit and existing infrastructure & Canonical evidence, provenance and reviewable comparison workflows & Source contracts, validation and release records & AI agent involvement recorded; individual component ownership unresolved \\
Revision and retrospective diagnostics & Preserved manuscript, sources and completed audits & Personal review and acceptance pending & Exact-model proof, coupled/numerical fixtures, outcome extraction and document checks & Codex produced the proof, diagnostics, new type analysis, comparator, benchmark, protocol and drafting; personal verification pending \\
\end{xltabular}
\endgroup
% END SOURCE BLOCK 151

% BEGIN SOURCE BLOCK 152 paragraph
Canonical representations preserve substantive fields while excluding volatile timestamps. Entry markers distinguish successive visits in padded slots; typed parsing separates input interpretation from deterministic metrics, and provenance rejects absent nodes \hyperref[source-s11]{[S11]}.
% END SOURCE BLOCK 152

% BEGIN SOURCE BLOCK 153 paragraph
Imported-bundle provenance reconciles signed success contributions with each run's own denominator, failing on mismatch rather than inventing matched identities. This generic software capability is distinct from the VEC joins; SUMO/TOS adapters do not implement it.
% END SOURCE BLOCK 153

% BEGIN SOURCE BLOCK 154 paragraph
Aggregate service-work and count carry preserve comparability and tractability without individual departures. A discrete-event redesign would require new model validation; faster scheduling cannot resolve that limitation.
% END SOURCE BLOCK 154

% BEGIN SOURCE BLOCK 155 paragraph
Original morning checks included matched inputs/actions, logit tolerance 0.00001, full/probe prefixes and separate workload reconstruction (maximum accumulated discrepancy about 0.00174 ms). The prior raw audit checked admissions and outcomes independently. These complementary checks were not repeated here and do not physically validate the model. Table~\ref{tab:12} links limitations to claims.
% END SOURCE BLOCK 155

% BEGIN SOURCE BLOCK 156 table
\begingroup
\small
\begin{xltabular}{\textwidth}{@{}YYY@{}}
\caption{Validity limits and their consequences for interpretation.}\label{tab:12} \\
\toprule
Limitation & Claim restricted & Consequence / evidence needed \\
\midrule\endfirsthead
\toprule
Limitation & Claim restricted & Consequence / evidence needed \\
\midrule\endhead
\bottomrule\endfoot
Batched arrivals, aggregate departures, backlog-only gate & Physical deadline feasibility & Interpret as evaluator outcomes; validate an event/transfer model against measurements \\
Global actual service work; immediate reservations & Distributed deployability & Measure estimation error and coordination delay before implementation claims \\
Four fleet draws, one actor/task seed per scenario & Population uncertainty and generalisation & Intervals condition on these controls; more independent draws and streams needed \\
Bundled scenario changes & Morning gain caused by lower density & Separate scenario analyses; factor-controlled evidence needed \\
Legacy off masks differ; E0/E1/E2 originals unavailable & Admission-consistent E1 attainment and gate-only E2b attribution & Historical score intervals remain; measurement impact is not assessable from available records \\
Exact arithmetic versus float32; finite fixtures and unobserved intermediate masks & Numerical prevalence and historical causal decomposition & Two-deadline exact bound proved; constructed rounding discrepancy and reduced adverse example do not establish fleet-frequency effects \\
Single-draw sensitivities & General staleness/forwarding robustness & Describe tested information and fixed-cost models only \\
Unconfirmed backup/access and personal ownership & Unrestricted reproduction and individual achievement & Portable compact checks supplied; raw access and component/AI allocation require confirmation \\
\end{xltabular}
\endgroup
% END SOURCE BLOCK 156

% BEGIN SOURCE BLOCK 157 paragraph
Portable verification regenerates central tables and accepts explicit September inputs. Original E0/E1/E2 raw outputs remain unavailable following author-reported deletion, with no known backup; other historical availability is unestablished. The completed September inventory/audits retain their recorded scope. Checksums and local presence are not off-machine preservation; an approved storage endpoint and examiner access remain unconfirmed \hyperref[source-s13]{[S13]}, \hyperref[source-s14]{[S14]}.
% END SOURCE BLOCK 157

% BEGIN SOURCE BLOCK 158 heading
\subsection{A sealed extension, not new confirmation evidence}\label{sec:3.7}
% END SOURCE BLOCK 158

% BEGIN SOURCE BLOCK 159 paragraph
The proposed study uses eight fresh joint fleet/evaluator seed pairs, (100,200) through (107,207), selected by an outcome-independent absence rule over recorded evaluation manifests. Four arms yield 32 planned morning evaluations; the delivered execution switch is false. Configuration generation, input/source identities, storage checks and runner dry-run are complete; full qualification and outcomes are unrun \hyperref[source-s15]{[S15]}.
% END SOURCE BLOCK 159

% BEGIN SOURCE BLOCK 160 paragraph
Evaluator seed changes observation descriptors, arrival counts, operational task types/sizes, separate observation/operational fading and service wobble. Fleet seed changes tier, EV status and initial SoC. All arms must match recorded exogenous streams within a block; observations, logits and actions are checked separately and allowed to diverge through legitimate feedback. The estimand is the total scheduling intervention under frozen weights, not fixed-action replay.
% END SOURCE BLOCK 160

% BEGIN SOURCE BLOCK 161 paragraph
Primary paired contrasts are per-task\ensuremath{-}ingress, common-target\ensuremath{-}ingress and per-task\ensuremath{-}round-robin. Bonferroni 95\% family intervals use eight independent joint blocks and t critical 3.128 (df=7), assuming approximately normal block effects. At an assumed SD of 1 percentage point, half-width is 1.106 points; a standardised effect of 1 has only approximately 44.7\% per-contrast power. Old fixed-task-seed variance cannot establish this design's precision. Failures stop the matrix and remain recorded; no outcome-driven retries, extra seeds or pooling with earlier draws are allowed. Workload awareness beyond spreading and robustness to fresh joint randomness remain open.
% END SOURCE BLOCK 161

% BEGIN SOURCE BLOCK 162 heading
\section{Conclusion}\label{sec:4}
% END SOURCE BLOCK 162

% BEGIN SOURCE BLOCK 163 paragraph
TrafficTwin establishes an implementation-dependent ranking reversal under a frozen vehicle policy: common-target placement underperformed ingress, while causal per-task placement exceeded it in both Manchester scenarios. The contribution is controlled evaluation and explanation of established mechanisms.
% END SOURCE BLOCK 163

% BEGIN SOURCE BLOCK 164 paragraph
RQ1 remains bounded by separate uncertainties. The archived E1 capacity-score contrast is statistically inconclusive under its historical implementation; interpretation as admission-consistent attainment retains the unquantified off-mask exception. Surviving summaries reproduce intervals, while unavailable original task records prevent resolving measurement impact. Equivalence is unsupported.
% END SOURCE BLOCK 164

% BEGIN SOURCE BLOCK 165 paragraph
For RQ2, gate-enabled ingress supplies the stronger comparator. E2b retains eligibility/scoring differences; later incident contrasts yield common-target \ensuremath{-}2.122 and per-task +0.527 percentage points. E2d is adaptive and reuses controls. For RQ3, the independent morning replication excludes its inspected pilot and reproduces the directions under simultaneous intervals: \ensuremath{-}3.232 and +3.899 points. Four fleet draws remain the replications. Descriptive accounting locates 284,821 gains and 12,842 losses; the two 100 ms task types supply almost all net gains, with no type-level net harm in these draws but some increased admitted misses.
% END SOURCE BLOCK 165

% BEGIN SOURCE BLOCK 166 paragraph
The exact-model proposition separates reconciliation from reselection. Two deadline values guarantee causal fixed-target admission within three simultaneous replacements in exact arithmetic. Gate/service/timing constraints explain clearing and repeated common destinations. Float32 boundary cases qualify executable equivalence without measuring historical prevalence; the reduced adverse example identifies why current workload need not maximise timely completions. These analyses do not uniquely decompose the full performance gap.
% END SOURCE BLOCK 166

% BEGIN SOURCE BLOCK 167 paragraph
RQ4 shows that identical 100 ms reports made that information treatment inactive, while live admission and reservations protect the tested stale-report model. Fixed forwarding cost reduced but did not eliminate the single-draw advantage through 10 ms; network congestion remains outside the model.
% END SOURCE BLOCK 167

% BEGIN SOURCE BLOCK 168 paragraph
The new host benchmark quantifies lower scheduler cost for causal scans on these fixtures and an additional reduction for cyclic targeting. Its companion comparator and sealed joint-randomness protocol make the next uncertainty explicit: whether workload awareness earns a deadline benefit over spreading when evaluator randomness also varies. That campaign remains unrun. Physical timing, generalisation, personal contribution, assessment-specific AI use and examiner access require their own evidence or author decisions; kernel speed and verification cannot substitute for them.
% END SOURCE BLOCK 168

% BEGIN SOURCE BLOCK 169 bibliography_open
\begin{thebibliography}{99}
% END SOURCE BLOCK 169

% BEGIN SOURCE BLOCK 170 bibliography_item
\bibitem{ref1} Y. Mao, C. You, J. Zhang, K. Huang and K. B. Letaief. ``A Survey on Mobile Edge Computing: The Communication Perspective.'' \emph{IEEE Communications Surveys \& Tutorials}, 19(4), 2322--2358, 2017. DOI: 10.1109/COMST.2017.2745201. \href{\detokenize{https://arxiv.org/abs/1701.01090}}{Author manuscript}.
% END SOURCE BLOCK 170

% BEGIN SOURCE BLOCK 171 bibliography_item
\bibitem{ref2} M. Harchol-Balter, M. E. Crovella and C. D. Murta. ``On Choosing a Task Assignment Policy for a Distributed Server System.'' \emph{Computer Performance Evaluation (TOOLS 1998)}, LNCS 1469, 231--242, 1998. DOI: 10.1007/3-540-68061-6\_19. \href{\detokenize{https://www.cs.cmu.edu/~harchol/Papers/tools.pdf}}{Author manuscript}.
% END SOURCE BLOCK 171

% BEGIN SOURCE BLOCK 172 bibliography_item
\bibitem{ref3} J. Schulman, F. Wolski, P. Dhariwal, A. Radford and O. Klimov. ``Proximal Policy Optimization Algorithms.'' arXiv:1707.06347, 2017. \href{\detokenize{https://arxiv.org/abs/1707.06347}}{Author manuscript}.
% END SOURCE BLOCK 172

% BEGIN SOURCE BLOCK 173 bibliography_item
\bibitem{ref4} C. Yu, A. Velu, E. Vinitsky, J. Gao, Y. Wang, A. Bayen and Y. Wu. ``The Surprising Effectiveness of PPO in Cooperative Multi-Agent Games.'' \emph{Advances in Neural Information Processing Systems}, 35, 24611--24624, 2022. \href{\detokenize{https://papers.nips.cc/paper/2022/file/9c1535a02f0ce079433344e14d910597-Paper-Datasets_and_Benchmarks.pdf}}{Proceedings paper}.
% END SOURCE BLOCK 173

% BEGIN SOURCE BLOCK 174 bibliography_item
\bibitem{ref5} P. Henderson, R. Islam, P. Bachman, J. Pineau, D. Precup and D. Meger. ``Deep Reinforcement Learning That Matters.'' \emph{Proceedings of the AAAI Conference on Artificial Intelligence}, 32(1), 2018. DOI: 10.1609/aaai.v32i1.11694. \href{\detokenize{https://ojs.aaai.org/index.php/AAAI/article/view/11694}}{Publisher record}.
% END SOURCE BLOCK 174

% BEGIN SOURCE BLOCK 175 bibliography_item
\bibitem{ref6} R. Agarwal, M. Schwarzer, P. S. Castro, A. Courville and M. G. Bellemare. ``Deep Reinforcement Learning at the Edge of the Statistical Precipice.'' \emph{Advances in Neural Information Processing Systems}, 34, 29304--29320, 2021. \href{\detokenize{https://proceedings.neurips.cc/paper/2021/hash/f514cec81cb148559cf475e7426eed5e-Abstract.html}}{Proceedings record}.
% END SOURCE BLOCK 175

% BEGIN SOURCE BLOCK 176 bibliography_item
\bibitem{ref7} R. Gorsane, O. Mahjoub, R. de Kock, R. Dubb, S. Singh and A. Pretorius. ``Towards a Standardised Performance Evaluation Protocol for Cooperative MARL.'' \emph{Advances in Neural Information Processing Systems}, 35, 5510--5521, 2022. \href{\detokenize{https://arxiv.org/abs/2209.10485}}{Author manuscript}.
% END SOURCE BLOCK 176

% BEGIN SOURCE BLOCK 177 bibliography_item
\bibitem{ref8} P. Alvarez Lopez, M. Behrisch, L. Bieker-Walz, J. Erdmann, Y.-P. Flötteröd, R. Hilbrich, L. Lücken, J. Rummel, P. Wagner and E. Wießner. ``Microscopic Traffic Simulation using SUMO.'' \emph{21st International Conference on Intelligent Transportation Systems (ITSC)}, 2575--2582, 2018. DOI: 10.1109/ITSC.2018.8569938. \href{\detokenize{https://elib.dlr.de/127994/}}{DLR author repository}.
% END SOURCE BLOCK 177

% BEGIN SOURCE BLOCK 178 bibliography_item
\bibitem{ref9} K. Ousterhout, P. Wendell, M. Zaharia and I. Stoica. ``Sparrow: Distributed, Low Latency Scheduling.'' \emph{24th ACM Symposium on Operating Systems Principles}, 69--84, 2013. DOI: 10.1145/2517349.2522716. \href{\detokenize{https://sigops.org/s/conferences/sosp/2013/papers/p69-ousterhout.pdf}}{Proceedings paper}.
% END SOURCE BLOCK 178

% BEGIN SOURCE BLOCK 179 bibliography_item
\bibitem{ref10} S. Vargaftik, I. Keslassy and A. Orda. ``LSQ: Load Balancing in Large-Scale Heterogeneous Systems With Multiple Dispatchers.'' \emph{IEEE/ACM Transactions on Networking}, 28(3), 1186--1198, 2020. DOI: 10.1109/TNET.2020.2980061. \href{\detokenize{https://webee.technion.ac.il/~isaac/p/ton20_lsq.pdf}}{Author manuscript}.
% END SOURCE BLOCK 179

% BEGIN SOURCE BLOCK 180 bibliography_item
\bibitem{ref11} J. Niño-Mora. ``Admission and routing of soft real-time jobs to multiclusters: Design and comparison of index policies.'' \emph{Computers \& Operations Research}, 39, 3431--3444, 2012. DOI: 10.1016/j.cor.2012.05.004. \href{\detokenize{https://arxiv.org/abs/2207.12815}}{Author manuscript, deposited 2022}.
% END SOURCE BLOCK 180

% BEGIN SOURCE BLOCK 181 bibliography_item
\bibitem{ref12} R. G. Sargent. ``Verification and Validation of Simulation Models.'' \emph{Proceedings of the 2011 Winter Simulation Conference}, 183--198, 2011. \href{\detokenize{https://www.informs-sim.org/wsc11papers/016.pdf}}{Proceedings paper}.
% END SOURCE BLOCK 181

% BEGIN SOURCE BLOCK 182 bibliography_item
\bibitem{ref13} L. Ying, R. Srikant and X. Kang. ``The Power of Slightly More than One Sample in Randomized Load Balancing.'' \emph{IEEE INFOCOM}, 1131--1139, 2015. \href{\detokenize{https://doi.org/10.1109/INFOCOM.2015.7218487}}{doi:10.1109/INFOCOM.2015.7218487}. \href{\detokenize{https://bpb-us-w2.wpmucdn.com/sites.coecis.cornell.edu/dist/9/287/files/2019/08/Srikant-1-YinSriKan14tech.pdf}}{Author technical manuscript}.
% END SOURCE BLOCK 182

% BEGIN SOURCE BLOCK 183 bibliography_item
\bibitem{ref14} JAX authors. ``Benchmarking JAX code.'' Official JAX documentation, accessed 8 September 2026. \href{\detokenize{https://docs.jax.dev/en/latest/benchmarking.html}}{Benchmarking guidance}. Used for timing procedure, not scheduler efficacy.
% END SOURCE BLOCK 183

% BEGIN SOURCE BLOCK 184 bibliography_close
\end{thebibliography}
% END SOURCE BLOCK 184

% BEGIN SOURCE BLOCK 185 appendix_open
\appendix
\renewcommand{\thetable}{\Alph{section}\arabic{table}}
% END SOURCE BLOCK 185

% BEGIN SOURCE BLOCK 186 heading
\appendixsection{app:A}{Evidence sources and reproducibility}
% END SOURCE BLOCK 186

% BEGIN SOURCE BLOCK 187 paragraph
The sources below are project evidence, separate from the scholarly references. Relative links resolve in the preserved repository checkout; the \href{\detokenize{https://github.com/Abdulla4akash/traffictwin/tree/04f3b6a95c7bc06c80ed95b54762f12861bba183}}{baseline tree} fixes their contents. The accompanying \href{\detokenize{../empirical_extension_2026-09-08/CLAIM_SOURCE_MAP.md}}{claim-to-source map} provides claim-level limits and additional receipts.
% END SOURCE BLOCK 187

% BEGIN SOURCE BLOCK 188 anchor
\phantomsection\label{source-s0}
% END SOURCE BLOCK 188

% BEGIN SOURCE BLOCK 189 paragraph
\textbf{S0 --- Accounting validity.} \href{\detokenize{../../evaluation/vec_followup_2026-09-07/historical_e0_e2d/experiments/E0/README.md}}{E0 report and evidence}. Full evidence commit \texttt{29d8945862b7df1bd5d7879ba1d980a388f6be0d}.
% END SOURCE BLOCK 189

% BEGIN SOURCE BLOCK 190 anchor
\phantomsection\label{source-s1}
% END SOURCE BLOCK 190

% BEGIN SOURCE BLOCK 191 paragraph
\textbf{S1 --- Waiting-room capacity.} \href{\detokenize{../../evaluation/vec_followup_2026-09-07/historical_e0_e2d/experiments/E1/README.md}}{E1 report}, with five-draw comparison and validation records.
% END SOURCE BLOCK 191

% BEGIN SOURCE BLOCK 192 anchor
\phantomsection\label{source-s2}
% END SOURCE BLOCK 192

% BEGIN SOURCE BLOCK 193 paragraph
\textbf{S2 --- Exploratory admission/placement decomposition.} \href{\detokenize{../../evaluation/vec_followup_2026-09-07/historical_e0_e2d/experiments/E2b/README.md}}{E2b report and factorial evidence}.
% END SOURCE BLOCK 193

% BEGIN SOURCE BLOCK 194 anchor
\phantomsection\label{source-s3}
% END SOURCE BLOCK 194

% BEGIN SOURCE BLOCK 195 paragraph
\textbf{S3 --- Incident replication and reversal.} \href{\detokenize{../../evaluation/vec_followup_2026-09-07/historical_e0_e2d/experiments/E2c/README.md}}{E2c report} and \href{\detokenize{../../evaluation/vec_followup_2026-09-07/historical_e0_e2d/experiments/E2d/README.md}}{E2d report}. Frozen E2d evaluator commit \texttt{2f63706f46319433a2ba3af1df97afd0e56a95d1}; exact E2c controls reused.
% END SOURCE BLOCK 195

% BEGIN SOURCE BLOCK 196 anchor
\phantomsection\label{source-s4}
% END SOURCE BLOCK 196

% BEGIN SOURCE BLOCK 197 paragraph
\textbf{S4 --- Experimental architecture and implementation.} \href{\detokenize{../../evaluation/vec_followup_2026-09-07/historical_e0_e2d/docs/METHODOLOGY.md}}{Frozen method}, \href{\detokenize{../../evaluation/vec_followup_2026-09-07/frozen_evaluator/eval/e2d_per_task_placement.py}}{per-task source}, \href{\detokenize{../../evaluation/vec_followup_2026-09-07/frozen_evaluator/eval/eval_sumo_stage1_mc.py}}{evaluator source}, \href{\detokenize{../empirical_extension_2026-09-08/experimental/vendor/vec_jax.py}}{environment source}, and \href{\detokenize{../../correspondence/sandra_randy_vec_progress_email_thread_2026-08-18.md}}{implementation correspondence}, PDF pages 1--2 for Randy's confirmations.
% END SOURCE BLOCK 197

% BEGIN SOURCE BLOCK 198 anchor
\phantomsection\label{source-s5}
% END SOURCE BLOCK 198

% BEGIN SOURCE BLOCK 199 paragraph
\textbf{S5 --- State-information sensitivity.} \href{\detokenize{../../evaluation/vec_followup_2026-09-07/frozen_evaluator/docs/RSU_STATE_DELAY.md}}{Timing contract}, \href{\detokenize{../../evaluation/vec_followup_2026-09-07/state-delay-pilot-2026-09-07/comparison.csv}}{five-cell comparison} and \href{\detokenize{../../evaluation/vec_followup_2026-09-07/state-delay-pilot-2026-09-07/mechanism_audit.json}}{report-state audit}.
% END SOURCE BLOCK 199

% BEGIN SOURCE BLOCK 200 anchor
\phantomsection\label{source-s6}
% END SOURCE BLOCK 200

% BEGIN SOURCE BLOCK 201 paragraph
\textbf{S6 --- Forwarding sensitivity.} \href{\detokenize{../../evaluation/vec_followup_2026-09-07/forwarding-sensitivity-2026-09-07/qualification.json}}{Qualification}, \href{\detokenize{../../evaluation/vec_followup_2026-09-07/forwarding-sensitivity-2026-09-07/forwarding_sensitivity.csv}}{outcomes} and \href{\detokenize{../../evaluation/vec_followup_2026-09-07/forwarding-sensitivity-2026-09-07/analysis_validation.json}}{analysis validation}.
% END SOURCE BLOCK 201

% BEGIN SOURCE BLOCK 202 anchor
\phantomsection\label{source-s7}
% END SOURCE BLOCK 202

% BEGIN SOURCE BLOCK 203 paragraph
\textbf{S7 --- Morning exploratory pilot.} \href{\detokenize{../../evaluation/vec_followup_2026-09-07/generalisation-pilot-2026-09-07/RESULTS.md}}{Results} and \href{\detokenize{../../evaluation/vec_followup_2026-09-07/generalisation-pilot-2026-09-07/input_validation.json}}{canonical input validation}.
% END SOURCE BLOCK 203

% BEGIN SOURCE BLOCK 204 anchor
\phantomsection\label{source-s8}
% END SOURCE BLOCK 204

% BEGIN SOURCE BLOCK 205 paragraph
\textbf{S8 --- Morning primary replication.} \href{\detokenize{../../evaluation/vec_followup_2026-09-07/generalisation-replication-2026-09-07/PROTOCOL.md}}{Protocol}, \href{\detokenize{../../evaluation/vec_followup_2026-09-07/generalisation-replication-2026-09-07/manifest.json}}{manifest}, \href{\detokenize{../../evaluation/vec_followup_2026-09-07/generalisation-replication-2026-09-07/paired_differences.csv}}{paired differences} and \href{\detokenize{../../evaluation/vec_followup_2026-09-07/generalisation-replication-2026-09-07/paired_intervals.csv}}{paired intervals}. Full follow-up evaluator commit \texttt{908bd10f86542de94fc38af90dd56c2ccc08cf9b}.
% END SOURCE BLOCK 205

% BEGIN SOURCE BLOCK 206 anchor
\phantomsection\label{source-s9}
% END SOURCE BLOCK 206

% BEGIN SOURCE BLOCK 207 paragraph
\textbf{S9 --- Five-RSU mechanism audit.} \href{\detokenize{../../evaluation/common_target_mechanism_audit_2026-09-07/evidence/REPORT.md}}{Completed post-hoc report} and \href{\detokenize{../../evaluation/common_target_mechanism_audit_2026-09-07/evidence/audit_results.json}}{audit results with input identities}.
% END SOURCE BLOCK 207

% BEGIN SOURCE BLOCK 208 anchor
\phantomsection\label{source-s10}
% END SOURCE BLOCK 208

% BEGIN SOURCE BLOCK 209 paragraph
\textbf{S10 --- Unrun falsification proposal.} \href{\detokenize{../../evaluation/common_target_mechanism_audit_2026-09-07/TIEBREAK_PREFIX_PROPOSAL.md}}{Rotating-tie-break prefix protocol}.
% END SOURCE BLOCK 209

% BEGIN SOURCE BLOCK 210 anchor
\phantomsection\label{source-s11}
% END SOURCE BLOCK 210

% BEGIN SOURCE BLOCK 211 paragraph
\textbf{S11 --- Artifact and preservation.} \href{\detokenize{../../architecture_current_release.md}}{Product architecture}, \href{\detokenize{../../quality/final_release_status_20260815.md}}{release/execution boundary}, \href{\detokenize{../../evaluation/vec_followup_2026-09-07/ARCHIVE_INVENTORY.json}}{evidence archive inventory} and \href{\detokenize{../../evaluation/vec_followup_2026-09-07/COMPACT_VERIFICATION.json}}{compact verification receipt}. The separate E3 Dynamic Resource V2 campaign remains unexecuted.
% END SOURCE BLOCK 211

% BEGIN SOURCE BLOCK 212 anchor
\phantomsection\label{source-s12}
% END SOURCE BLOCK 212

% BEGIN SOURCE BLOCK 213 paragraph
\textbf{S12 --- Contribution and repair history.} Read-only inspection of \texttt{vec\_env} commits \href{\detokenize{https://github.com/Abdulla4akash/vec_env/commit/4cb7c06aab1b455d68f71ea9143172463784dea9}}{4cb7c06}, \href{\detokenize{https://github.com/Abdulla4akash/vec_env/commit/0f01f4d2082d3e8b735e74a873095ab8eeba37cc}}{0f01f4d} and \href{\detokenize{https://github.com/Abdulla4akash/vec_env/commit/2f63706f46319433a2ba3af1df97afd0e56a95d1}}{2f63706}, and TrafficTwin \href{\detokenize{https://github.com/Abdulla4akash/traffictwin/blob/29d8945862b7df1bd5d7879ba1d980a388f6be0d/tests/test_validate_e0_smoke.py}}{E0 validation tests at 29d894}. Source comments explicitly credit Abdulla's questions; commit authorship does not establish independent human authorship.
% END SOURCE BLOCK 213

% BEGIN SOURCE BLOCK 214 anchor
\phantomsection\label{source-s13}
% END SOURCE BLOCK 214

% BEGIN SOURCE BLOCK 215 paragraph
\textbf{S13 --- Retrospective gap closure, 8 September 2026.} The \href{\detokenize{../gap_closure_2026-09-08/verification/README.md}}{verification entry point} documents the field contract, source bindings, 87-file authentication, 19-run audit, \href{\detokenize{../gap_closure_2026-09-08/verification/results/OUTCOME_TRANSITIONS.json}}{23 paired transitions}, \href{\detokenize{../gap_closure_2026-09-08/verification/results/KERNEL_RESULTS.json}}{440-case reference} and production comparison. These completed checks and their historical sources remain unchanged; they are not rerun here.
% END SOURCE BLOCK 215

% BEGIN SOURCE BLOCK 216 anchor
\phantomsection\label{source-s14}
% END SOURCE BLOCK 216

% BEGIN SOURCE BLOCK 217 paragraph
\textbf{S14 --- Production-domain mechanism analysis.} \href{\detokenize{../production_mechanism_2026-09-08/analysis/MATHEMATICS.md}}{Model and proof}, \href{\detokenize{../production_mechanism_2026-09-08/analysis/PROTOCOL.md}}{initial diagnostic protocol}, \href{\detokenize{../production_mechanism_2026-09-08/analysis/FOLLOWUP_PROTOCOL.md}}{bounded follow-up protocol}, \href{\detokenize{../production_mechanism_2026-09-08/analysis/diagnostics.py}}{executable diagnostics}, \href{\detokenize{../production_mechanism_2026-09-08/analysis/results/DIAGNOSTICS.json}}{full results}, \href{\detokenize{../production_mechanism_2026-09-08/analysis/results/NUMERICAL_FOLLOWUP.json}}{numerical follow-up}, \href{\detokenize{../production_mechanism_2026-09-08/analysis/results/ADVERSE_PRODUCTION.json}}{adverse example and helper checks}, \href{\detokenize{../production_mechanism_2026-09-08/analysis/results/DIMENSION_CHECK.json}}{dimension checks} and \href{\detokenize{../production_mechanism_2026-09-08/analysis/results/OUTCOME_ACCOUNTING.json}}{outcome accounting}. This is retrospective analysis added after the completed experiments and gap closure.
% END SOURCE BLOCK 217

% BEGIN SOURCE BLOCK 218 anchor
\phantomsection\label{source-s15}
% END SOURCE BLOCK 218

% BEGIN SOURCE BLOCK 219 paragraph
\textbf{S15 --- Retrospective empirical/engineering extension.} \href{\detokenize{../empirical_extension_2026-09-08/analysis/PROTOCOL.md}}{Declared analysis protocol}, \href{\detokenize{../empirical_extension_2026-09-08/analysis/results/TASK_TYPES.json}}{task-type counts and bindings}, \href{\detokenize{../empirical_extension_2026-09-08/analysis/benchmark.py}}{benchmark implementation} and \href{\detokenize{../empirical_extension_2026-09-08/analysis/results/BENCHMARK.json}}{measured results}, \href{\detokenize{../empirical_extension_2026-09-08/experimental/round_robin.py}}{cyclic implementation}, \href{\detokenize{../empirical_extension_2026-09-08/analysis/results/COMPATIBILITY_TESTS.json}}{compatibility evidence}, \href{\detokenize{../empirical_extension_2026-09-08/confirmation/PROTOCOL.md}}{sealed confirmation protocol}, \href{\detokenize{../empirical_extension_2026-09-08/confirmation/DRY_RUN.json}}{dry-run receipt}, and \href{\detokenize{../empirical_extension_2026-09-08/README.md}}{verification route}. New campaign status: not executed.
% END SOURCE BLOCK 219

% BEGIN SOURCE BLOCK 220 paragraph
The runtime is CPython 3.11.15, JAX/JAXlib 0.4.30 and NumPy 1.26.4, CPU, x64 disabled. Historical manifests retain their own identities. Full reproduction also needs the frozen actor, traces and authorised inputs. E0/E1/E2 originals remain unavailable following author-reported deletion, with no known backup; other historical raw availability is unestablished.
% END SOURCE BLOCK 220

% BEGIN SOURCE BLOCK 221 paragraph
\textbf{Assistance and author review.} Codex assisted substantive research, drafting, audits, proof, diagnostics and this extension's analysis, implementation, tests, benchmark and protocol. These are AI-assisted project artifacts. The \href{\detokenize{../empirical_extension_2026-09-08/AUTHOR_INPUTS.md}}{author checklist} separates recorded contributions from pending personal verification, assessment-specific AI authority and approval; no completed student review is inferred.
% END SOURCE BLOCK 221

% BEGIN SOURCE BLOCK 222 heading
\appendixsection{app:B}{Audit coverage and diagnostic boundaries}
% END SOURCE BLOCK 222

% BEGIN SOURCE BLOCK 223 paragraph
Table~\ref{tab:B1} retains the completed gap-closure audit coverage. Those raw checks and joins were not repeated for this revision.
% END SOURCE BLOCK 223

% BEGIN SOURCE BLOCK 224 table
\begingroup
\small
\begin{xltabular}{\textwidth}{@{}YYYY@{}}
\caption{Current availability and fresh audit coverage, 8 September 2026. Historical receipts remain unchanged.}\label{tab:B1} \\
\toprule
Records & Available evidence / coverage & Completed gap-closure checks & Consequence \\
\midrule\endfirsthead
\toprule
Records & Available evidence / coverage & Completed gap-closure checks & Consequence \\
\midrule\endhead
\bottomrule\endfoot
E0: two ten-step repeats and full seed 0 & Surviving receipts, manifests, sources; full receipt records 834,120 cap rejections & Not assessable from available records & Original raw outputs unavailable; deletion reported by author; no known backup \\
E1: five seeds \ensuremath{\times} three caps & Fifteen run summaries/receipts; seed-0 middle cap reuses E0 & Not assessable from available records & Original interval reproduced from summaries; scoring impact unresolved \\
E2 off/jsq/dla seed 0 and E2b reused cells & Hash-bound summaries and receipts; E2b reuses E2 exactly & Not assessable from available records & No new cell or historical task join inferred from reuse \\
E2b new ingress / E2c / E2d originals & Compact sources/records survive; original raw availability not established & Not assessable from available records & No deletion or fresh raw-validation claim extended to these studies \\
September morning primary & Twelve full runs, seeds 0/2/3/4; 1,744,116 offered per cell & Zero discrepancies in each inspected run & Supports these gate-enabled records; excludes pilot from inference \\
September morning pilot & Two full seed-1 runs, 1,744,116 offered each & Zero discrepancies in each inspected run & Remains exploratory, outside primary inference \\
September incident sensitivity pilot & Five full seed-1 runs, 13,076,234 offered each & Zero discrepancies in each inspected run & Does not replace historical E2d records \\
Other September arrays & Remaining inventory entries authenticated only & Integrity-only in the prior audit & Qualification/preflight files are not additional fleet replications \\
\end{xltabular}
\endgroup
% END SOURCE BLOCK 224

% BEGIN SOURCE BLOCK 225 paragraph
The per-run JSON records preserve offered/admitted counts, explicit rejection and unavailability categories, false-success and latency discrepancy counts, penalty-equal ambiguity, summary reconciliation and source bindings. A rejected task's finite latency is not necessarily a false success. For the 19 available full runs, both false-success and non-penalty rejection-latency counts are zero, and fixed-history admission-consistent scores equal their archived values. Historical unavailable cells are labelled ``not assessable,'' never zero. Receipt checks are reported as past observations, not freshly rerun validation.
% END SOURCE BLOCK 225

% BEGIN SOURCE BLOCK 226 paragraph
Table~\ref{tab:B2} details primary morning ingress\ensuremath{\to}per-task task accounting. ``Both'' columns partition the same offered population; no task-level confidence intervals are constructed.
% END SOURCE BLOCK 226

% BEGIN SOURCE BLOCK 227 table
\begingroup
\small
\begin{xltabular}{\textwidth}{@{}YYYYYY@{}}
\caption{Descriptive task transitions from ingress to per-task placement; four primary morning draws, pilot excluded.}\label{tab:B2} \\
\toprule
Fleet seed & Failure \ensuremath{\to} success & Success \ensuremath{\to} failure & Success both & Failure both & Net successes \\
\midrule\endfirsthead
\toprule
Fleet seed & Failure \ensuremath{\to} success & Success \ensuremath{\to} failure & Success both & Failure both & Net successes \\
\midrule\endhead
\bottomrule\endfoot
0 & 71,152 & 3,340 & 1,545,724 & 123,900 & 67,812 \\
2 & 83,543 & 4,640 & 1,522,085 & 133,848 & 78,903 \\
3 & 70,862 & 2,913 & 1,550,789 & 119,552 & 67,949 \\
4 & 59,264 & 1,949 & 1,569,719 & 113,184 & 57,315 \\
\end{xltabular}
\endgroup
% END SOURCE BLOCK 227

% BEGIN SOURCE BLOCK 228 paragraph
All 23 eligible within-study pairs are retained in the machine-readable record. Task identity uses matched trace-row/substep/slot coordinates and verified stream identities; sampled sizes are source-derived correspondence, not independently observed fields. Later queue divergence prevents these descriptive transitions from identifying a unique causal contribution of individual placements.
% END SOURCE BLOCK 228

% BEGIN SOURCE BLOCK 229 paragraph
The scalar suite includes 440 primary cases and ten deletion trials, below its 500-case limit. The five-task case is minimal under single-candidate deletion, not a globally minimal counterexample over every numeric parameter. A's retained mask remains unchanged by the diagnostic code: the fourth replacement is inspected only. C's repeated complete passes restart from the same state. No full-evaluator tie-break prefix, actor inference, SUMO execution, campaign or retraining was performed.
% END SOURCE BLOCK 229

% BEGIN SOURCE BLOCK 230 heading
\appendixsection{app:C}{Exact-model result and numerical boundary}
% END SOURCE BLOCK 230

% BEGIN SOURCE BLOCK 231 paragraph
The full proof and executable listings are in S14. Its admission map uses fixed eligibility $E$, targets, initial workload/count, nonnegative service, exclusive earlier-candidate offsets and strict backlog/capacity tests. The causal scan is its unique fixed point. From arbitrary masks, a destination with $n$ eligible candidates settles within $n$ replacements; from $E$, the stronger bound is $\min(n,2d-1)$ for $d$ distinct deadlines. Nonincreasing deadlines require only one replacement. Empty destinations need none; zero work still consumes count. These deductions concern the stated model and do not imply scheduler optimality.
% END SOURCE BLOCK 231

% BEGIN SOURCE BLOCK 232 paragraph
The initial 671 inputs comprise 128 general mathematical fixtures, 324 coupled fixtures, 210 deliberate threshold fixtures and nine completion boundaries. Bounded adverse reductions/transfers and numerical follow-ups bring the total to 719, within the predeclared ceiling of 800. Protocol amendments precede their specific follow-ups; no outcome-driven expansion beyond that budget occurred. Full negative and unchanged results remain in the machine-readable record.
% END SOURCE BLOCK 232

% BEGIN SOURCE BLOCK 233 paragraph
The numerical alternating case uses seven type3 services of 2.5252525806427 ms followed by one type1 service of 21.11111068725586 ms, all with 100 ms deadlines. Initial $W=82.32322692871094$ ms and $Q=4$, $C=6{,}220$. Four earlier type1 services of $W/4$ construct this queue in a separate local admission substep; their noise multiplier lies within 0.9--1.1. This establishes local kernel compatibility under explicit arrays, not exact RNG/actor/traffic reachability.
% END SOURCE BLOCK 233

% BEGIN SOURCE BLOCK 234 paragraph
Exact arithmetic admits all eight, with the last gate margin about 0.000005007 ms; the last task nevertheless misses completion. The instrumented float32 mirror alternates $11111111\to11111110\to11111111\to11111110\to11111111$. The unchanged production helper confirms the retained third mask, recomputed offsets and capacity-failure count. Its causal float32 scan also rejects the last candidate, though via its own rounded gate. Thus this fixture distinguishes exact from implemented arithmetic and exposes an unstable label, not a demonstrated A\ensuremath{\to}B success improvement. The mirror's intermediate masks are constructed diagnostics, not recorded September fields.
% END SOURCE BLOCK 234

% BEGIN SOURCE BLOCK 235 paragraph
Four padding checks retain the alternating mask at $N=215/R=9$ and $N=2{,}488/R=10$, with the active block at either end. Their target is imposed: a busy target would not be the common argmin while others are empty. A final fixture prepends the four services, starts every queue empty and selects RSU0 at the correct $N=215/R=9$ dimensions. Its third mask is stable, although the last admission differs from exact arithmetic. This comparison prevents treating fixed-target instability as demonstrated on a full common-target trajectory.
% END SOURCE BLOCK 235

% BEGIN SOURCE BLOCK 236 paragraph
One separate inclusive-completion boundary also differs: saved-offset arithmetic rounds a just-late exact completion to an on-time float32 value. This is a zero-transfer calculation, not execution of the full latency scorer. Neither numerical example measures historical frequency or changes the reported fleet estimates.
% END SOURCE BLOCK 236

% BEGIN SOURCE BLOCK 237 paragraph
The two-RSU adverse case is minimal under the tested shared-position deletions and one-substep reduction. It retains production task/deadline/capacity values but changes infrastructure dimensions; its twelve transfers to the studied RSU/substep parameters retain no adverse result. No fixture frequency is interpreted as a fleet probability, and no new confidence intervals, raw audits or full-system runs are reported.
% END SOURCE BLOCK 237

% BEGIN SOURCE BLOCK 238 heading
\appendixsection{app:D}{Type aggregation and engineering evidence}
% END SOURCE BLOCK 238

% BEGIN SOURCE BLOCK 239 paragraph
This new retrospective analysis reads only eight primary-morning task arrays, authenticates their hashes, and aggregates categories by type. The prior 19-run admission validation and 23 task joins are reused, not repeated. All-task totals and paired net changes reconcile. The following table exposes every type/draw/arm; attainment uses the offered column. Other terminal failures are available in the linked CSV and unchanged within each pair.
% END SOURCE BLOCK 239

% BEGIN SOURCE BLOCK 240 table
\begingroup
\small
\begin{xltabular}{\textwidth}{@{}YYYYYYYY@{}}
\caption{Type outcomes. I = ingress, P = per-task; offered and category counts are tasks, attainment is percent. Gate rejection and admitted misses are disjoint.}\label{tab:D1} \\
\toprule
Seed / type & Arm & Offered & Success & Attainment & Admitted & Gate rejected & Admitted miss \\
\midrule\endfirsthead
\toprule
Seed / type & Arm & Offered & Success & Attainment & Admitted & Gate rejected & Admitted miss \\
\midrule\endhead
\bottomrule\endfoot
0/1 & I & 349,045 & 252,283 & 72.278 & 328,923 & 19,548 & 76,640 \\
0/1 & P & 349,045 & 272,106 & 77.957 & 347,605 & 866 & 75,499 \\
0/2 & I & 523,060 & 516,034 & 98.657 & 522,134 & 46 & 6,100 \\
0/2 & P & 523,060 & 516,106 & 98.671 & 522,180 & 0 & 6,074 \\
0/3 & I & 872,011 & 780,747 & 89.534 & 821,868 & 48,572 & 41,121 \\
0/3 & P & 872,011 & 828,664 & 95.029 & 868,284 & 2,156 & 39,620 \\
2/1 & I & 349,045 & 242,629 & 69.512 & 324,165 & 24,149 & 81,536 \\
2/1 & P & 349,045 & 261,794 & 75.003 & 347,317 & 997 & 85,523 \\
2/2 & I & 523,060 & 515,409 & 98.537 & 521,901 & 40 & 6,492 \\
2/2 & P & 523,060 & 515,471 & 98.549 & 521,941 & 0 & 6,470 \\
2/3 & I & 872,011 & 768,687 & 88.151 & 809,915 & 60,362 & 41,228 \\
2/3 & P & 872,011 & 828,363 & 94.995 & 867,771 & 2,506 & 39,408 \\
3/1 & I & 349,045 & 253,921 & 72.747 & 330,364 & 18,383 & 76,443 \\
3/1 & P & 349,045 & 274,260 & 78.574 & 348,537 & 210 & 74,277 \\
3/2 & I & 523,060 & 517,080 & 98.857 & 522,533 & 1 & 5,453 \\
3/2 & P & 523,060 & 517,086 & 98.858 & 522,534 & 0 & 5,448 \\
3/3 & I & 872,011 & 782,701 & 89.758 & 824,760 & 46,357 & 42,059 \\
3/3 & P & 872,011 & 830,305 & 95.217 & 870,539 & 578 & 40,234 \\
4/1 & I & 349,045 & 260,449 & 74.618 & 333,800 & 14,596 & 73,351 \\
4/1 & P & 349,045 & 280,703 & 80.420 & 348,250 & 146 & 67,547 \\
4/2 & I & 523,060 & 516,064 & 98.662 & 521,928 & 1 & 5,864 \\
4/2 & P & 523,060 & 516,069 & 98.663 & 521,929 & 0 & 5,860 \\
4/3 & I & 872,011 & 795,155 & 91.186 & 834,388 & 35,739 & 39,233 \\
4/3 & P & 872,011 & 832,211 & 95.436 & 869,818 & 309 & 37,607 \\
\end{xltabular}
\endgroup
% END SOURCE BLOCK 240

% BEGIN SOURCE BLOCK 241 paragraph
The \href{\detokenize{../empirical_extension_2026-09-08/analysis/results/BENCHMARK.json}}{benchmark record} retains lowering/compilation times, median, quartiles, p95, minima/maxima, compiler buffer estimates and exact source/runtime identities for all 16 configurations. \href{\detokenize{../empirical_extension_2026-09-08/analysis/results/BENCHMARK_TIMES.csv}}{Raw call durations} permit descriptive timing checks. \href{\detokenize{../empirical_extension_2026-09-08/confirmation/PRECISION.json}}{Precision sensitivity} concerns assumed future block variance; it is not a result of the unrun confirmation study.
% END SOURCE BLOCK 241

\end{document}
